% Luc Destombe --- licence CC BY-NC-SA 4.0
% https://creativecommons.org/licenses/by-nc-sa/4.0/deed.fr
% Fichier autonome : compiler avec pdflatex, deux fois (signets, nombre de pages).
\documentclass[11pt,a4paper]{article}
\makeatletter
\newif\iflycee@unecolonne
\lycee@unecolonnetrue
\newif\iflycee@palatino
\lycee@palatinotrue

\RequirePackage[utf8]{inputenc}
\RequirePackage[T1]{fontenc}
\RequirePackage[french]{babel}
\RequirePackage{etoolbox}

\newcommand{\ShorthandsOff}{\@ifundefined{shorthandoff}{}{\shorthandoff{;:!?}}}
\newcommand{\ShorthandsOn}{\@ifundefined{shorthandon}{}{\shorthandon{;:!?}}}
\ShorthandsOff
\AtBeginDocument{\ShorthandsOn}
\AtBeginEnvironment{tikzpicture}{\ShorthandsOff}

\RequirePackage{geometry}
\RequirePackage{fancyhdr}
\RequirePackage{lastpage}
\RequirePackage{multicol}
\RequirePackage{array}
\RequirePackage{enumitem}
\RequirePackage{xcolor}
\RequirePackage{needspace}

\RequirePackage{amsmath,amssymb}
\RequirePackage{mathpazo}
\DeclareMathSizes{10.95}{10.95}{8}{5.5}
\begingroup\catcode`\_=\active \catcode`\^=\active
\gdef_#1{\sb{\scriptscriptstyle #1}}
\gdef^#1{\sp{\scriptscriptstyle\lycee@moinsepais #1}}
\endgroup
\newcommand\lycee@moins{\mathbin{%
  \setbox\z@\hbox{$\m@th\scriptscriptstyle\lycee@moinsorig$}%
  \dimen@=.55\wd\z@ \advance\dimen@-.5pt
  \hbox to.75\wd\z@{\hss$\m@th\scriptscriptstyle\vcenter{\hbox{%
    \tikz\draw[line width=.5pt,line cap=round](0,0)--(\the\dimen@,0);}}$\hss}}}
\begingroup\lccode`\~=`\- \lowercase{\endgroup\let~}\lycee@moins
\newcommand\lycee@moinsepais{\mathcode`\-="8000 }
\AtBeginDocument{\mathchardef\lycee@moinsorig=\mathcode`\-\relax}
\AtBeginDocument{\catcode`\_=12 \mathcode`\_="8000
  \catcode`\^=12 \mathcode`\^="8000 }
\newcommand\lycee@exposants{%
  \fontdimen13\textfont2=.48\fontdimen6\textfont2
  \fontdimen14\textfont2=.48\fontdimen6\textfont2
  \fontdimen15\textfont2=.48\fontdimen6\textfont2
  \fontdimen13\scriptfont2=.48\fontdimen6\scriptfont2
  \fontdimen14\scriptfont2=.48\fontdimen6\scriptfont2
  \fontdimen15\scriptfont2=.48\fontdimen6\scriptfont2}
\every@math@size\expandafter{\the\every@math@size\lycee@exposants}
\newlength{\retraitcalculs}
\setlength{\retraitcalculs}{2em}

\RequirePackage{tikz}
\usetikzlibrary{arrows.meta,calc,babel,shapes.misc}
\RequirePackage[most]{tcolorbox}
\tcbuselibrary{skins,breakable}

\definecolor{coulDef}{HTML}{1F4E79}
\definecolor{coulProp}{HTML}{7030A0}
\definecolor{coulRem}{HTML}{C55A11}
\definecolor{coulEx}{HTML}{2E7D32}
\definecolor{coulExo}{HTML}{B03A2E}
\definecolor{coulPart}{HTML}{1F4E79}
\definecolor{coulMeth}{HTML}{1B6E4F}
\definecolor{coulCourbe}{HTML}{0B5ED7}
\definecolor{coulCourbeB}{HTML}{C00000}
\definecolor{coulCourbeC}{HTML}{2E7D32}
\definecolor{coulGrille}{HTML}{8C8C8C}
\definecolor{coulRep}{HTML}{1B6E4F}
\definecolor{coulBar}{HTML}{2E7D32}

\newcommand{\R}{\mathbb{R}}
\newcommand{\Rs}{\mathbb{R}^{*}}
\renewcommand{\le}{\leqslant}
\renewcommand{\ge}{\geqslant}
\renewcommand{\approx}{\simeq}
\newcommand{\vect}[1]{\overrightarrow{#1}}
\newcommand{\norme}[1]{\left\lVert #1 \right\rVert}
\newcommand{\intervalle}[2]{\left[#1\,;#2\right]}
\RequirePackage{cancel}
\renewcommand{\CancelColor}{\color{coulExo}}
\newcommand{\fois}[1]{\times{\color{coulExo}#1}}
\newcommand{\barre}[1]{\cancel{#1}}

\tikzset{grille/.style={coulGrille,line width=0.4pt}}
\newcommand{\quadrillage}[4]{\draw[grille,step=1] (#1,#2) grid (#3,#4);}
\tikzset{
  croix/.style={cross out,draw,line width=0.6pt,inner sep=0pt,minimum size=4pt},
  croixdroite/.style={croix,rotate=45,line width=0.8pt},
}
\newcommand{\pt}[3]{\path (#1) node[croix]{} node[#2,font=\small,inner sep=2.5pt]{$#3$};}

\newcommand{\lycee@repere}[4]{%
  \draw[grille,step=1] (#1,#2) grid (#3,#4);
  \draw[-{Stealth[length=2mm]},thick] (#1,0)--({#3+0.4},0) node[below left=-1pt]{$x$};
  \draw[-{Stealth[length=2mm]},thick] (0,#2)--(0,{#4+0.4}) node[below left=-1pt]{$y$};
  \node[below left,font=\scriptsize,inner sep=1.5pt] at (0,0) {$0$};}
\newcommand{\lycee@gradx}[1]{\foreach \k/\l in {#1}{%
  \draw (\k,0.07)--(\k,-0.07) node[below,font=\scriptsize,inner sep=1.5pt]{$\l$};}}
\newcommand{\lycee@grady}[1]{\foreach \k/\l in {#1}{%
  \draw (0.07,\k)--(-0.07,\k) node[left,font=\scriptsize,inner sep=1.5pt]{$\l$};}}
\newcommand{\lycee@intff}[2]{\left[#1\,;#2\right]}
\newcommand{\lycee@intfo}[2]{\left[#1\,;#2\right[}
\newcommand{\lycee@intof}[2]{\left]#1\,;#2\right]}
\newcommand{\lycee@intoo}[2]{\left]#1\,;#2\right[}
\newcommand{\lycee@ens}[1]{\left\{#1\right\}}
\AtBeginDocument{%
  \providecommand{\repere}{\lycee@repere}%
  \providecommand{\gradx}{\lycee@gradx}%
  \providecommand{\grady}{\lycee@grady}%
  \providecommand{\Cf}{\mathcal{C}\sb{\scriptscriptstyle f}}%
  \providecommand{\Cg}{\mathcal{C}\sb{\scriptscriptstyle g}}%
  \providecommand{\intff}{\lycee@intff}%
  \providecommand{\intfo}{\lycee@intfo}%
  \providecommand{\intof}{\lycee@intof}%
  \providecommand{\intoo}{\lycee@intoo}%
  \providecommand{\ens}{\lycee@ens}}

\newlist{questions}{enumerate}{3}
\setlist[questions,1]{label=\arabic*\textdegree),leftmargin=*,itemsep=2pt,topsep=3pt}
\setlist[questions,2]{label=\alph*),leftmargin=*,itemsep=1pt,topsep=2pt}
\setlist[questions,3]{label=\roman*),leftmargin=*,itemsep=1pt,topsep=1pt}
\setlist[itemize]{label=\textbullet,leftmargin=*,itemsep=2pt,topsep=3pt}

\newcommand{\besoinplace}[1]{\par\penalty\@M\begingroup
  \dimen@=#1\relax\dimen@ii=\pagegoal\advance\dimen@ii-\pagetotal
  \ifdim\dimen@>\dimen@ii\ifdim\dimen@ii>\z@\vfil\fi\break\fi\endgroup}

\newcommand{\lycee@pied}{\thepage/\pageref{LastPage}}
\newcommand{\entete}[2]{%
  \pagestyle{fancy}\fancyhf{}%
  \fancyhead[L]{\footnotesize\color{coulPart}\textbf{#1}}%
  \fancyhead[R]{\footnotesize\color{coulPart}\textbf{#2}}%
  \fancyfoot[C]{\footnotesize\lycee@pied}%
  \renewcommand{\headrulewidth}{0.4pt}%
  \renewcommand{\headrule}{\hbox to\headwidth{%
    \color{coulPart!40}\leaders\hrule height \headrulewidth\hfill}}}

\RequirePackage{ccicons}
\newcommand{\lycee@licence}{{\scriptsize\color{gray}%
  \copyright\ Luc Destombe --- \ccbyncsa\ CC BY-NC-SA 4.0}}
\newif\iflycee@licence \lycee@licencetrue
\newcommand{\sanslicence}{\lycee@licencefalse}
\AtBeginDocument{\iflycee@licence\AddToHookNext{shipout/foreground}{%
  \put(0,0){\rlap{\kern\dimexpr1in+\hoffset+\oddsidemargin+\textwidth\relax
    \llap{\raisebox{-\dimexpr1in+\voffset+\topmargin+\headheight+\headsep
      +\textheight+\footskip\relax}{\lycee@licence}}}}}\fi}

\newcommand{\formulesserrees}{%
  \setlength{\abovedisplayskip}{3pt}\setlength{\belowdisplayskip}{3pt}%
  \setlength{\abovedisplayshortskip}{2pt}\setlength{\belowdisplayshortskip}{3pt}}

\setlength{\parindent}{0pt}

\RequirePackage[implicit=false,hidelinks,pdfencoding=auto,psdextra,bookmarksopen,
  bookmarksopenlevel=1,pdfcreator={},pdfproducer={}]{hyperref}
\RequirePackage{bookmark}
\hypersetup{pdfauthor={Luc Destombe},pdfkeywords={CC BY-NC-SA 4.0}}
\newcounter{lycee@signet}
\newcommand{\signet}[3][]{\stepcounter{lycee@signet}%
  \edef\lycee@dest{\if\relax\detokenize{#1}\relax
    signet.\arabic{lycee@signet}\else#1\fi}%
  \hypertarget{\lycee@dest}{}\bookmark[level=#2,dest=\lycee@dest]{#3}}
\pdfstringdefDisableCommands{%
  \def\R{R}\def\vect#1{#1}\def\Cf{Cf}\def\Cg{Cg}\def\fois#1{×#1}%
  \def\barre#1{#1}\def\intervalle#1#2{[#1;#2]}\def\intff#1#2{[#1;#2]}%
  \def\textsuperscript#1{#1}\def\dfrac#1#2{#1/#2}\def\frac#1#2{#1/#2}%
  \def\vec#1{#1⃗}}%
\RequirePackage[official]{eurosym}

\tcbset{exercice corrige/.style={
  enhanced,breakable,before upper=\raggedright,
  colback=white,colframe=coulExo,
  boxrule=0.8pt,arc=2pt,
  left=6pt,right=6pt,top=8pt,bottom=5pt,
  before skip=9pt,after skip=4pt,
  coltitle=white,fonttitle=\bfseries\small,
  attach boxed title to top left={xshift=8pt,yshift=-\tcboxedtitleheight/2},
  boxed title style={colback=coulExo,arc=2pt,boxrule=0pt}}}
\newcounter{exo}
\newtcolorbox{exercice}[1][]{exercice corrige,
  phantom={\signet[exercice-\arabic{exo}]{2}{Exercice \arabic{exo}}},
  title={Exercice \theexo\ifx\relax#1\relax\else{} --- #1\fi}}
\newenvironment{exo}[1][\relax]{\refstepcounter{exo}\begin{exercice}[#1]}{\end{exercice}}

\geometry{top=1.8cm,bottom=1cm,left=1cm,right=1cm,headheight=14pt,footskip=0.6cm}
\renewcommand{\lycee@pied}{\thepage}

\renewcommand{\theexo}{\ifnum\value{exo}<10 0\fi\arabic{exo}}
\tcbset{exercice corrige/.style={
  enhanced,breakable,before upper=\raggedright,
  colback=white,colframe=coulExo,
  boxrule=0.9pt,arc=2pt,
  left=8pt,right=8pt,top=8pt,bottom=6pt,
  before skip=8pt,after skip=6pt,
  coltitle=white,fonttitle=\bfseries,
  attach boxed title to top left={xshift=8pt,yshift=-\tcboxedtitleheight/2},
  boxed title style={colback=coulExo,arc=2pt,boxrule=0pt}}}

\newcommand{\entetefiche}[2]{%
  \begin{center}
    \begin{tcolorbox}[enhanced,width=0.92\linewidth,colback=white,colframe=coulPart,
        boxrule=1.4pt,arc=3pt,halign=center,top=6pt,bottom=6pt]
      {\large\bfseries\color{coulPart} CORRIGÉ --- FICHE D'EXERCICES}\\[3pt]
      {\normalsize\bfseries #1}\\[2pt]
      {\normalsize #2}
    \end{tcolorbox}
  \end{center}}

\newcounter{partie}
\newcommand{\partie}[1]{%
  \stepcounter{partie}%
  \besoinplace{8\baselineskip}\vspace{4pt}%
  \begin{tcolorbox}[enhanced,colback=coulPart!8,colframe=coulPart,boxrule=0pt,
      leftrule=3pt,arc=0pt,left=6pt,right=4pt,top=3pt,bottom=3pt,
      before skip=6pt,after skip=2pt]
    \bfseries\color{coulPart}\Roman{partie} --- #1\signet{1}{\Roman{partie} --- #1}
  \end{tcolorbox}}
\newcommand{\partiesansnum}[1]{%
  \besoinplace{8\baselineskip}\vspace{4pt}%
  \begin{tcolorbox}[enhanced,colback=coulPart!8,colframe=coulPart,boxrule=0pt,
      leftrule=3pt,arc=0pt,left=6pt,right=4pt,top=3pt,bottom=3pt,
      before skip=6pt,after skip=2pt]
    \bfseries\color{coulPart}#1\signet{1}{#1}
  \end{tcolorbox}}

\setlist[questions,1]{label=\arabic*\textdegree),leftmargin=*,itemsep=3pt,topsep=3pt}
\setlist[questions,2]{label=\alph*),leftmargin=*,itemsep=2pt,topsep=2pt}
\setlist[questions,3]{label=\roman*),leftmargin=*,itemsep=1pt,topsep=1pt}
\newlist{expr}{enumerate}{1}
\setlist[expr]{label={\Alph*\ $=$},leftmargin=*,itemsep=2pt,topsep=3pt}

\newcommand{\res}[1]{{\color{coulRep}#1}}
\newcommand{\rem}[1]{\par\smallskip{\small\itshape\color{black!60}$\rightarrow$ #1}}
\newenvironment{calculs}{\par\smallskip\noindent$\displaystyle\begin{aligned}[t]}%
  {\end{aligned}$\par\smallskip}
\newcommand{\justif}[1]{\quad\text{\small\itshape\color{black!60}#1}}
\newcommand{\cas}[1]{\par\smallskip\noindent\textbf{\color{coulExo}#1}\nobreak}
\newcommand{\explique}[2]{\par\smallskip\noindent
  \begin{minipage}[t]{0.6\linewidth}\raggedright #1\end{minipage}\hfill
  \begin{minipage}[t]{0.37\linewidth}\raggedright\leavevmode
    {\small\itshape\color{black!60}#2}\end{minipage}\par}

\setlength{\parskip}{2pt}

\makeatother
\usepackage{tkz-tab}
\entete{Chapitre 5 --- Variations, signe et inéquations --- corrigé}{Seconde}
\setlength{\parindent}{0pt}
\setlength{\jot}{5pt}
\let\casligne\cas
\renewcommand{\cas}[1]{\casligne{#1}\par\nopagebreak}
\renewcommand{\res}[1]{\leavevmode{\color{coulRep}#1}}

\definecolor{coulInt}{HTML}{D62828}
\definecolor{coulIntB}{HTML}{1565C0}
\definecolor{coulIntC}{HTML}{2E7D32}
\NewDocumentCommand{\morceau}{O{coulInt} m m}{%
  \draw[#1,line width=2.2pt,opacity=0.85] (#2,0)--(#3,0);}
\NewDocumentCommand{\crochet}{O{coulInt} m m}{%
  \ifx[#3%
    \draw[#1,line width=1.2pt] ($(#2,0)+(2.6pt,5pt)$) -- ++(-2.6pt,0) -- ++(0,-10pt) -- ++(2.6pt,0);
  \else
    \draw[#1,line width=1.2pt] ($(#2,0)+(-2.6pt,5pt)$) -- ++(2.6pt,0) -- ++(0,-10pt) -- ++(-2.6pt,0);
  \fi}
\NewDocumentCommand{\solx}{O{coulInt} m}{\path[#1] (#2,0) node[croix]{};}

\newenvironment{lecture}{\par\noindent\begin{minipage}[t]{0.55\linewidth}\vspace{0pt}\raggedright}%
  {\end{minipage}\hfill}
\newenvironment{figlecture}{\begin{minipage}[t]{0.43\linewidth}\vspace{0pt}\centering}%
  {\end{minipage}\par\smallskip}
\newcommand{\gris}[1]{\leavevmode{\small\itshape\color{black!60}#1}}

\newcommand{\signeaffine}[4]{%
  \begin{tikzpicture}[scale=0.85]
    \tkzTabInit[lgt=2.4,espcl=2.2]{$x$/1,#1/1}{$-\infty$,#2,$+\infty$}
    \tkzTabLine{,#3,z,#4,}
  \end{tikzpicture}}

\newcommand{\courbeI}{%
  \repere{-4}{-3}{6}{4}
  \gradx{-3,-2,-1,1,2,3,4,5} \grady{-2,-1,1,2,3}
  \draw[coulCourbe,line width=1.1pt]
    (-3,-2) .. controls (-2,0.5) and (-1,3) .. (0,3)
            .. controls (1,3) and (2,-1) .. (3,-1)
            .. controls (3.67,-1) and (4.33,1) .. (5,2);
  \path[coulCourbe] (-3,-2) node[croix]{} (5,2) node[croix]{};
  \node[coulCourbe,font=\footnotesize] at (-2.6,1) {$\Cf$};}
\newcommand{\courbeXXVII}{%
  \repere{-5}{-3}{6}{3}
  \gradx{-4,-3,-2,-1,1,2,3,4,5} \grady{-2,-1,1,2}
  \draw[coulCourbe,line width=1.1pt]
    (-4,2) .. controls (-3.67,1.33) and (-3.33,0.44) .. (-3,0)
           .. controls (-2.33,-0.89) and (-1.67,-2) .. (-1,-2)
           .. controls (-0.33,-2) and (0.33,-0.89) .. (1,0)
           .. controls (1.33,0.44) and (1.67,2) .. (2,2)
           .. controls (2.67,2) and (3.33,0.89) .. (4,0)
           .. controls (4.33,-0.44) and (4.67,-1.33) .. (5,-2);
  \path[coulCourbe] (-4,2) node[croix]{} (5,-2) node[croix]{};
  \node[coulCourbe,font=\footnotesize] at (-4.4,1) {$\Cf$};}
\newcommand{\droitesXX}{%
  \repere{-4}{-4}{5}{5}
  \gradx{-3,-2,-1,1,2,3,4} \grady{-3,-2,-1,1,2,3,4}
  \draw[coulCourbe,line width=1.1pt] (-1,-4)--(3.5,5);
  \draw[coulCourbeB,line width=1.1pt] (-4,4)--(5,-0.5);
  \draw[coulCourbeC,line width=1.1pt] (-4,3)--(3,-4);
  \draw[black!70,line width=1.1pt] (-4,-2.667)--(5,3.333);
  \node[coulCourbe,font=\footnotesize,right] at (3.4,4.6) {$(d_1)$};
  \node[coulCourbeB,font=\footnotesize,above] at (-3.4,4.2) {$(d_2)$};
  \node[coulCourbeC,font=\footnotesize,left] at (-3.9,2.6) {$(d_3)$};
  \node[black!70,font=\footnotesize,right] at (5,3.3) {$(d_4)$};}

\begin{document}

\entetefiche{Chapitre 5 --- Variations, signe et inéquations}{Seconde}

\partie{Fonction croissante, fonction décroissante}

\begin{exo}[Lire les variations]
\cas{1)}
\explique{$f$ est définie sur $\res{\intff{-3}{5}}$.}{la courbe commence en $x=-3$ et
s'arrête en $x=5$ (les deux croix) : on lit les abscisses de ses extrémités}

\cas{2)}
\begin{lecture}
On parcourt la courbe \textbf{de gauche à droite}, comme un profil d'étape de montagne.

\smallskip
De $x=-3$ à $x=0$, la courbe monte (de $-2$ à $3$) : $f$ est \res{croissante sur
$\intff{-3}{0}$}.

\smallskip
De $x=0$ à $x=3$, elle descend (de $3$ à $-1$) : $f$ est \res{décroissante sur
$\intff{0}{3}$}.

\smallskip
De $x=3$ à $x=5$, elle remonte (de $-1$ à $2$) : $f$ est \res{croissante sur
$\intff{3}{5}$}.

\smallskip
\gris{Sur la figure, chaque intervalle est tracé en couleur sur l'axe des abscisses : on
descend des points où la courbe change de sens jusqu'à l'axe, et on lit $0$ et $3$.}
\end{lecture}
\begin{figlecture}
\begin{tikzpicture}[x=0.42cm,y=0.42cm]
  \courbeI
  \draw[black!50,dashed] (0,3)--(0,0) (3,-1)--(3,0);
  \morceau{-3}{0}\morceau[coulIntB]{0}{3}\morceau[coulIntC]{3}{5}
\end{tikzpicture}
\end{figlecture}
\rem{Erreur fréquente : écrire les intervalles avec les hauteurs ($\intff{-2}{3}$, par
exemple). Les variations s'écrivent avec des valeurs de $x$, lues sur l'axe des
abscisses.}
\rem{En $0$ et en $3$, la courbe change de sens : ces valeurs sont la borne de deux
intervalles voisins, avec des crochets fermés.}
\end{exo}

\begin{exo}[Une randonnée]
\cas{1)}
\explique{$f$ est \res{croissante sur $\intff{8}{11}$}.}{de 8~h à 11~h, Inès monte : son
altitude augmente}
\explique{$f$ est \res{constante sur $\intff{11}{12}$}.}{de 11~h à 12~h, elle se repose
au refuge : l'altitude ne change pas}
\explique{$f$ est \res{décroissante sur $\intff{12}{15}$}.}{de 12~h à 15~h, elle
redescend : son altitude diminue}

\cas{2)}
\explique{$9$ et $10$ sont dans $\intff{8}{11}$, où $f$ est croissante.}{on cherche un
intervalle qui contient les deux nombres et où $f$ varie dans un seul sens}
\explique{$9<10$, donc $\res{f(9)\le f(10)}$ : à 10~h, Inès est plus haut qu'à 9~h.}{une
fonction croissante conserve l'ordre}
\explique{$13$ et $14$ sont dans $\intff{12}{15}$, où $f$ est décroissante.}{même
démarche}
\explique{$13<14$, donc $\res{f(13)\ge f(14)}$ : à 14~h, elle est plus bas qu'à
13~h.}{une fonction décroissante renverse l'ordre}
\rem{On ne connaît aucune altitude, sauf celle du parking : on compare sans calculer, avec
le sens de variation seulement.}
\end{exo}

\begin{exo}[Comparer des images]
\cas{1)}
\explique{On part de $2<5$.}{on compare d'abord les nombres de départ}
\explique{a) $g$ croissante : $\res{g(2)\le g(5)}$.}{l'ordre est conservé}
\explique{b) $g$ décroissante : $\res{g(2)\ge g(5)}$.}{l'ordre est renversé}
\explique{c) $g$ constante : $\res{g(2)=g(5)}$.}{$g(x)$ garde toujours la même valeur}

\cas{2)}
\explique{$-4$ et $-1$ sont dans $\intff{-5}{5}$, et $-4<-1$.}{sur la droite graduée,
$-4$ est à gauche de $-1$}
\explique{$f$ est croissante sur cet intervalle, donc $\res{f(-4)\le f(-1)}$.}{l'ordre est
conservé}
\rem{Erreur fréquente : croire $-4>-1$ parce que $4>1$. Avec des nombres négatifs, on pense
à la droite graduée.}
\end{exo}

\begin{exo}[Conserver ou renverser l'ordre]
\cas{1)}
\explique{$1<3<4<8$, et ces quatre nombres sont dans $\intff{0}{10}$.}{on range d'abord
les nombres de départ}
\explique{$f$ est croissante, donc $\res{f(1)\le f(3)\le f(4)\le f(8)}$.}{les images
sont rangées dans le même ordre}

\cas{2)}
\explique{$g$ est décroissante : les images sont rangées dans l'ordre inverse,
$g(1)\ge g(3)\ge g(4)\ge g(8)$.}{plus $x$ est grand, plus $g(x)$ est petit}
\explique{Dans l'ordre croissant : $\res{g(8)\le g(4)\le g(3)\le g(1)}$.}{on relit la ligne
précédente de droite à gauche}

\cas{3)}
\explique{Une fonction croissante \res{conserve} l'ordre ; une fonction décroissante
\res{renverse} l'ordre.}{c'est ce qu'on a vu en 1) et en 2)}

\cas{4)}
\explique{\res{Non}. Supposons $a>b$, c'est-à-dire $b<a$.}{on regarde ce qui se passerait}
\explique{$f$ est croissante, donc $f(b)\le f(a)$.}{l'ordre est conservé}
\explique{Or on sait que $f(a)<f(b)$ : c'est impossible. Donc $a$ n'est pas plus grand que
$b$.}{les deux inégalités se contredisent}
\rem{Si $a=b$, on aurait $f(a)=f(b)$ : ce n'est pas possible non plus. On a donc forcément
$a<b$.}
\end{exo}

\begin{exo}[Tracer une courbe]
\begin{lecture}
On place d'abord $A(0\,;1)$ et $B(6\,;3)$, puis on relie ces points en respectant les
variations.

\smallskip
1) Courbe bleue : elle monte de $A$ à $B$, sans jamais redescendre.

\smallskip
\gris{elle peut monter plus ou moins vite : il y a bien d'autres réponses}

\smallskip
2) Courbe rouge : elle descend de $A$ jusqu'à un point d'abscisse $2$, par exemple
$(2\,;-1)$, puis monte jusqu'à $B$.

\smallskip
\gris{Le point d'abscisse $2$ est plus bas que $A$ : $f$ décroît de $0$ à $2$, donc
$f(2)\le f(0)$, soit $f(2)\le 1$. La courbe y a une tangente horizontale : c'est là
qu'elle change de sens.}
\end{lecture}
\begin{figlecture}
\begin{tikzpicture}[x=0.5cm,y=0.5cm]
  \repere{-1}{-2}{7}{4}
  \gradx{1,2,3,4,5,6} \grady{-1,1,2,3}
  \draw[coulCourbe,line width=1pt]
    (0,1) .. controls (2,1.6) and (4,2.2) .. (6,3);
  \draw[coulCourbeB,line width=1pt]
    (0,1) .. controls (0.67,0.2) and (1.33,-1) .. (2,-1)
          .. controls (3.33,-1) and (4.67,1.6) .. (6,3);
  \pt{0,1}{above left}{A} \pt{6,3}{above left}{B}
\end{tikzpicture}
\end{figlecture}
\rem{Une courbe de fonction ne revient jamais en arrière : à chaque $x$ correspond un seul
point.}
\end{exo}

\begin{exo}[Vrai ou faux]
\cas{1)}
\explique{\res{Vrai}.}{$\intff{2}{5}$ est contenu dans $\intff{0}{10}$}
\explique{Si la courbe monte sur tout $\intff{0}{10}$, elle monte aussi sur le morceau qui
va de $x=2$ à $x=5$.}{pour $a<b$ dans $\intff{2}{5}$, $a$ et $b$ sont aussi dans
$\intff{0}{10}$, donc $f(a)\le f(b)$}

\cas{2)}
\begin{lecture}
\res{Faux}. Un seul contre-exemple suffit : une courbe qui monte, redescend, puis remonte.

\smallskip
Ci-contre : $f(1)=0$, $f(2)=3$, $f(3)=-1$ et $f(4)=1$.

\smallskip
On a bien $f(1)<f(4)$, car $0<1$.

\smallskip
Mais $2<3$ et $f(2)>f(3)$ : $f$ n'est pas croissante sur $\intff{1}{4}$.

\smallskip
\gris{comparer les images des deux bornes ne dit rien de ce qui se passe entre elles}
\end{lecture}
\begin{figlecture}
\begin{tikzpicture}[x=0.5cm,y=0.5cm]
  \repere{-1}{-2}{5}{4}
  \gradx{1,2,3,4} \grady{-1,1,2,3}
  \draw[coulCourbe,line width=1pt]
    (1,0) .. controls (1.33,1) and (1.67,3) .. (2,3)
          .. controls (2.33,3) and (2.67,-1) .. (3,-1)
          .. controls (3.33,-1) and (3.67,0.33) .. (4,1);
  \path[coulCourbe] (1,0) node[croix]{} (2,3) node[croix]{} (3,-1) node[croix]{}
    (4,1) node[croix]{};
\end{tikzpicture}
\end{figlecture}

\cas{3)}
\explique{\res{Faux}. Contre-exemple : la fonction $f$ de l'exercice 1, sur
$\intff{0}{5}$.}{elle descend de $x=0$ à $x=3$, puis remonte}
\explique{Elle n'est pas croissante sur $\intff{0}{5}$ : $0<3$ et $f(0)>f(3)$ ($3>-1$). Elle
n'est pas décroissante non plus : $3<5$ et $f(3)<f(5)$ ($-1<2$).}{une fonction peut n'être
ni croissante ni décroissante sur un intervalle}
\rem{Pour montrer qu'une affirmation est fausse, un seul contre-exemple suffit ; pour
montrer qu'elle est vraie, un exemple ne suffit pas.}
\end{exo}

\partie{Maximum, minimum}

\begin{exo}[Lire des extremums]
\begin{lecture}
\textbf{\color{coulExo}1)} Sur $\intff{-3}{5}$, on regarde toute la courbe.

\smallskip
Le point le plus haut est $(0\,;3)$ : le maximum est \res{$3$, atteint en $0$}.

\smallskip
Le point le plus bas est $(-3\,;-2)$ : le minimum est \res{$-2$, atteint en $-3$}.

\smallskip
\gris{Il est au bord : $-2$ est plus bas que le creux $(3\,;-1)$.}

\medskip
\textbf{\color{coulExo}2)} Sur $\intff{0}{5}$, on ne regarde que la partie de la courbe
entre $x=0$ et $x=5$.

\smallskip
Maximum : \res{$3$, atteint en $0$} ; minimum : \res{$-1$, atteint en $3$}.

\smallskip
\gris{le point $(-3\,;-2)$ n'est plus dans cette partie}

\medskip
\textbf{\color{coulExo}3)} Sur $\intff{3}{5}$, la courbe monte de $(3\,;-1)$ à $(5\,;2)$.

\smallskip
Minimum : \res{$-1$, atteint en $3$} ; maximum : \res{$2$, atteint en $5$}.

\smallskip
\gris{une fonction croissante sur un intervalle a son minimum à gauche et son maximum à
droite}
\end{lecture}
\begin{figlecture}
\begin{tikzpicture}[x=0.42cm,y=0.42cm]
  \courbeI
  \path[coulCourbe] (0,3) node[croix]{} (3,-1) node[croix]{};
  \draw[black!50,dashed] (3,-1)--(3,0) (5,2)--(5,0);
\end{tikzpicture}
\end{figlecture}
\rem{Le maximum est une image (une ordonnée) ; on dit en plus pour quelle valeur de $x$ il
est atteint. Sur un intervalle plus petit, l'extremum peut changer, et il peut être au
bord.}
\end{exo}

\begin{exo}[Tracer à partir d'une description]
\cas{1)}
\begin{lecture}
On traduit chaque donnée par un point de la courbe :

\smallskip
$f(-2)=0$ : point $(-2\,;0)$ ;\par
maximum $4$ atteint en $1$ : point $(1\,;4)$ ;\par
minimum $-3$ atteint en $4$ : point $(4\,;-3)$ ;\par
$f(6)=2$ : point $(6\,;2)$.

\smallskip
\gris{$f(a)=b$ : point d'abscisse $a$ et d'ordonnée $b$}

\smallskip
Puis on relie les points : la courbe monte de $(-2\,;0)$ à $(1\,;4)$, descend jusqu'à
$(4\,;-3)$, puis remonte jusqu'à $(6\,;2)$.

\smallskip
\gris{Tangente horizontale au sommet et au creux ; la courbe ne dépasse jamais $4$ et ne
descend jamais sous $-3$.}
\end{lecture}
\begin{figlecture}
\begin{tikzpicture}[x=0.4cm,y=0.4cm]
  \repere{-3}{-4}{7}{5}
  \gradx{-2,-1,1,2,3,4,5,6} \grady{-3,-2,-1,1,2,3,4}
  \draw[coulCourbe,line width=1pt]
    (-2,0) .. controls (-1,2) and (0,4) .. (1,4)
           .. controls (2,4) and (3,-3) .. (4,-3)
           .. controls (4.67,-3) and (5.33,0.33) .. (6,2);
  \path[coulCourbe] (-2,0) node[croix]{} (1,4) node[croix]{} (4,-3) node[croix]{}
    (6,2) node[croix]{};
\end{tikzpicture}
\end{figlecture}

\cas{2)}
\explique{Sur $\intff{4}{6}$, $f$ est croissante, de $f(4)=-3$ à $f(6)=2$.}{on ne regarde
que ce morceau de courbe}
\explique{Minimum : \res{$-3$, atteint en $4$} ; maximum : \res{$2$, atteint en
$6$}.}{croissante : minimum à gauche, maximum à droite}
\rem{Le maximum de $f$ sur $\intff{-2}{6}$ est $4$, mais sur $\intff{4}{6}$, c'est $2$ :
l'intervalle compte.}
\end{exo}

\begin{exo}[Les marées]
\cas{1)}
\explique{$f$ est \res{décroissante sur $\intff{2}{8}$}.}{la mer descend de la pleine mer
($12$~m à 2~h) à la basse mer ($2$~m à 8~h)}
\explique{$f$ est \res{croissante sur $\intff{8}{14}$}.}{elle remonte jusqu'à $13$~m à
14~h}
\explique{$f$ est \res{décroissante sur $\intff{14}{20}$}.}{elle redescend jusqu'à $1$~m
à 20~h}

\cas{2)}
\explique{Sur $\intff{2}{20}$, le maximum est \res{$13$~m, atteint à 14~h}.}{plus grande
des hauteurs connues : $12$ et $13$ aux pleines mers}
\explique{Le minimum est \res{$1$~m, atteint à 20~h}.}{plus petite des hauteurs aux
basses mers : $2$ et $1$}
\explique{Sur $\intff{2}{8}$, la mer ne fait que descendre : maximum \res{$12$~m, à
2~h} ; minimum \res{$2$~m, à 8~h}.}{décroissante : maximum à gauche, minimum à droite}

\cas{3)}
\explique{On place les quatre points $(2\,;12)$, $(8\,;2)$, $(14\,;13)$ et $(20\,;1)$, puis
on les relie en respectant les variations.}{tangente horizontale aux pleines et basses
mers : la mer y change de sens}
\begin{center}
\begin{tikzpicture}[x=0.2cm,y=0.4cm]
  \draw[grille,xstep=2,ystep=1] (0,0) grid (20,14);
  \draw[-{Stealth[length=2mm]},thick] (0,0)--(21.5,0) node[right,font=\footnotesize]{$t$ (h)};
  \draw[-{Stealth[length=2mm]},thick] (0,0)--(0,15.5) node[right,font=\footnotesize]{hauteur (m)};
  \foreach \k in {2,4,...,20} \draw (\k,0.2)--(\k,-0.2) node[below,font=\scriptsize,inner sep=1.5pt]{$\k$};
  \foreach \k in {2,4,...,14} \draw (0.8,\k)--(-0.8,\k) node[left,font=\scriptsize,inner sep=1.5pt]{$\k$};
  \node[below left,font=\scriptsize,inner sep=1.5pt] at (0,0) {$0$};
  \draw[coulCourbe,line width=1pt]
    (2,12) .. controls (4,12) and (6,2) .. (8,2)
           .. controls (10,2) and (12,13) .. (14,13)
           .. controls (16,13) and (18,1) .. (20,1);
  \path[coulCourbe] (2,12) node[croix]{} (8,2) node[croix]{} (14,13) node[croix]{}
    (20,1) node[croix]{};
\end{tikzpicture}
\end{center}
\rem{Un carreau vaut $2$~h en abscisse : le point de 8~h est au 4\textsuperscript{e}
carreau. Erreur fréquente : compter un carreau pour une heure.}
\end{exo}

\partie{Tableau de variations}

\begin{exo}[De la courbe au tableau]
\explique{Première ligne : $-3$ ; $0$ ; $3$ ; $5$.}{bornes $-3$ et $5$ ; changements de sens
en $0$ et en $3$}
\explique{Seconde ligne : une flèche qui monte, une qui descend, une qui monte.}{croissante,
décroissante, croissante (exercice 1)}
\explique{Au bout des flèches : $f(-3)=-2$, $f(0)=3$, $f(3)=-1$ et $f(5)=2$.}{ordonnées des
quatre croix, lues sur la courbe}

\smallskip
\noindent\begin{minipage}[c]{0.5\linewidth}\centering
\begin{tikzpicture}[scale=0.85]
  \tkzTabInit[lgt=1.4,espcl=2,deltacl=0.6]{$x$/1,$f(x)$/2}{$-3$,$0$,$3$,$5$}
  \tkzTabVar{-/$-2$,+/$3$,-/$-1$,+/$2$}
\end{tikzpicture}
\end{minipage}\hfill
\begin{minipage}[c]{0.45\linewidth}\centering
\begin{tikzpicture}[x=0.42cm,y=0.42cm]
  \courbeI
  \path[coulCourbe] (0,3) node[croix]{} (3,-1) node[croix]{};
\end{tikzpicture}
\end{minipage}
\par\smallskip
\rem{Sur la première ligne, seulement des valeurs de $x$ ; les hauteurs vont au bout des
flèches. Un nombre en haut est un sommet de la courbe, un nombre en bas un creux (ou une
extrémité).}
\end{exo}

\begin{exo}[Du tableau à la courbe]
\cas{1)}
\explique{$g$ est définie sur $\res{\intff{-4}{5}}$.}{première et dernière valeurs de la
première ligne}

\cas{2)}
\explique{$g$ est \res{décroissante sur $\intff{-4}{-1}$}, de $2$ à $-3$.}{la première
flèche descend}
\explique{$g$ est \res{croissante sur $\intff{-1}{2}$}, de $-3$ à $3$.}{la deuxième
flèche monte}
\explique{$g$ est \res{décroissante sur $\intff{2}{5}$}, de $3$ à $-1$.}{la troisième
flèche descend}

\cas{3)}
\begin{lecture}
On lit un point au bout de chaque flèche : $(-4\,;2)$, $(-1\,;-3)$, $(2\,;3)$ et
$(5\,;-1)$.

\smallskip
\gris{valeur de $x$ sur la première ligne, image au bout de la flèche}

\smallskip
On place ces quatre points, puis on les relie : la courbe descend, monte, puis
redescend, avec une tangente horizontale en $(-1\,;-3)$ et en $(2\,;3)$.

\smallskip
\gris{une réponse possible parmi d'autres : le tableau ne dit pas comment la courbe
s'incurve entre deux points}
\end{lecture}
\begin{figlecture}
\begin{tikzpicture}[x=0.4cm,y=0.4cm]
  \repere{-5}{-4}{6}{4}
  \gradx{-4,-3,-2,-1,1,2,3,4,5} \grady{-3,-2,-1,1,2,3}
  \draw[coulCourbe,line width=1pt]
    (-4,2) .. controls (-3,0.5) and (-2,-3) .. (-1,-3)
           .. controls (0,-3) and (1,3) .. (2,3)
           .. controls (3,3) and (4,0.3) .. (5,-1);
  \path[coulCourbe] (-4,2) node[croix]{} (-1,-3) node[croix]{} (2,3) node[croix]{}
    (5,-1) node[croix]{};
\end{tikzpicture}
\end{figlecture}
\end{exo}

\begin{exo}[Exploiter un tableau]
\cas{1)}
\explique{$-4$ et $-2$ sont dans $\intff{-5}{-1}$, où $k$ est croissante.}{on cherche
l'intervalle qui contient les deux nombres}
\explique{$-4<-2$, donc $\res{k(-4)\le k(-2)}$.}{croissante : l'ordre est conservé}
\explique{$1$ et $3$ sont dans $\intff{-1}{4}$, où $k$ est décroissante.}{même
intervalle}
\explique{$1<3$, donc $\res{k(1)\ge k(3)}$.}{décroissante : l'ordre est renversé}

\cas{2)}
\explique{$-3$ est dans $\intff{-5}{-1}$ et $2$ dans $\intff{-1}{4}$ : les deux nombres ne
sont pas dans un même intervalle où $k$ varie dans un seul sens.}{le sens de variation ne
s'applique qu'à deux nombres d'un même intervalle}
\explique{Sur $\intff{-5}{-1}$, $k$ monte de $-2$ à $4$ : $k(-3)$ est entre $-2$ et $4$.
Sur $\intff{-1}{4}$, $k$ descend de $4$ à $0$ : $k(2)$ est entre $0$ et $4$.}{on encadre
chaque image avec le tableau}
\explique{Les deux images peuvent valoir $1$ et $3$, ou $3$ et $1$ : \res{on ne peut pas
les comparer}.}{deux courbes respectant le tableau peuvent donner les deux réponses}

\cas{3)}
\explique{Le maximum de $k$ est \res{$4$, atteint en $-1$}.}{plus grand nombre de la
seconde ligne}
\explique{Le minimum de $k$ est \res{$-2$, atteint en $-5$}.}{on compare les deux nombres du
bas : $-2<0$}
\rem{Erreur fréquente : prendre pour minimum le dernier nombre, $0$. On compare tous les
nombres du bas des flèches.}
\end{exo}

\begin{exo}[Extremums et tableau]
\cas{1)}
\explique{Sur $\intff{-4}{6}$ : maximum \res{$5$, atteint en $0$} ; minimum \res{$-3$,
atteint en $3$}.}{on compare tous les nombres de la seconde ligne : $5>1$ et $-3<2$}
\explique{Sur $\intff{-4}{0}$ : maximum \res{$5$, atteint en $0$} ; minimum \res{$2$,
atteint en $-4$}.}{on ne garde que la première flèche, qui monte de $2$ à $5$}
\explique{Sur $\intff{3}{6}$ : minimum \res{$-3$, atteint en $3$} ; maximum \res{$1$,
atteint en $6$}.}{on ne garde que la dernière flèche, qui monte de $-3$ à $1$}

\cas{2)}
\explique{Pour tout $x$ de $\intff{-4}{6}$ : $\res{-3\le h(x)\le 5}$, donc $\res{m=-3}$ et
$\res{M=5}$ conviennent.}{le minimum et le maximum de $h$ sur $\intff{-4}{6}$}
\rem{Tout nombre $m$ inférieur à $-3$, et tout nombre $M$ supérieur à $5$, conviennent
aussi ($m=-10$ et $M=8$, par exemple) ; le minimum et le maximum donnent l'encadrement le
plus précis.}
\end{exo}

\begin{exo}[D'une description au tableau]
\cas{1)}
\explique{Première ligne : $-3$ ; $1$ ; $4$ ; $6$.}{bornes $-3$ et $6$ ; changements de sens
en $1$ et en $4$}
\explique{Seconde ligne : flèches qui descendent, montent, puis descendent, avec les images
données au bout.}{$f(-3)=5$, $f(1)=-2$, $f(4)=3$ et $f(6)=0$}
\begin{center}
\begin{tikzpicture}[scale=0.85]
  \tkzTabInit[lgt=1.4,espcl=2,deltacl=0.6]{$x$/1,$f(x)$/2}{$-3$,$1$,$4$,$6$}
  \tkzTabVar{+/$5$,-/$-2$,+/$3$,-/$0$}
\end{tikzpicture}
\end{center}

\cas{2)}
\explique{Le maximum de $f$ est \res{$5$, atteint en $-3$}.}{on compare les nombres du
haut : $5>3$ ; le maximum est au bord}
\explique{Le minimum de $f$ est \res{$-2$, atteint en $1$}.}{on compare les nombres du bas :
$-2<0$}

\cas{3)}
\begin{lecture}
On place les points $(-3\,;5)$, $(1\,;-2)$, $(4\,;3)$ et $(6\,;0)$, puis on les relie : la
courbe descend, monte, puis redescend.

\smallskip
\gris{tangente horizontale au creux $(1\,;-2)$ et au sommet $(4\,;3)$}
\end{lecture}
\begin{figlecture}
\begin{tikzpicture}[x=0.36cm,y=0.36cm]
  \repere{-4}{-3}{7}{6}
  \gradx{-3,-2,-1,1,2,3,4,5,6} \grady{-2,-1,1,2,3,4,5}
  \draw[coulCourbe,line width=1pt]
    (-3,5) .. controls (-1.67,1.67) and (-0.33,-2) .. (1,-2)
           .. controls (2,-2) and (3,3) .. (4,3)
           .. controls (4.67,3) and (5.33,1) .. (6,0);
  \path[coulCourbe] (-3,5) node[croix]{} (1,-2) node[croix]{} (4,3) node[croix]{}
    (6,0) node[croix]{};
\end{tikzpicture}
\end{figlecture}
\end{exo}

\partie{Fonctions affines et linéaires}

\begin{exo}[Reconnaître une fonction affine]
\explique{Pour chaque fonction, on cherche la forme $ax+b$ : $a$ est le nombre qui multiplie
$x$, $b$ le nombre seul.}{méthode 5 du cours ; linéaire si $b=0$, constante si $a=0$}

\cas{$f$}
\explique{$f(x)=5x+2$ est affine : $\res{a=5}$ et $\res{b=2}$.}{déjà sous la forme $ax+b$}

\cas{$g$}
\begin{calculs}
g(x) &= 6-2x &&\justif{on recopie l'expression}\\
g(x) &= -2x+6 &&\justif{on écrit d'abord le terme en $x$}
\end{calculs}
\explique{$g$ est affine : $\res{a=-2}$ et $\res{b=6}$.}{$a$ n'est pas le premier nombre
écrit, c'est celui qui multiplie $x$}

\cas{$h$}
\begin{calculs}
h(x) &= \frac{x}{3} &&\justif{on recopie l'expression}\\
h(x) &= \frac{1}{3}x &&\justif{diviser par $3$, c'est multiplier par $\frac13$}
\end{calculs}
\explique{$h$ est affine, et même \res{linéaire} : $\res{a=\frac{1}{3}}$ et
$\res{b=0}$.}{pas de nombre seul : $b=0$}

\cas{$k$}
\explique{$k(x)=-4$ est affine, et même \res{constante} : $\res{a=0}$ et
$\res{b=-4}$.}{$k(x)=0x-4$ : pas de terme en $x$}

\cas{$m$}
\explique{$m(x)=x^{2}+3$ \res{n'est pas affine}.}{il y a un terme en $x^{2}$, qui ne
s'écrit pas $ax$}

\cas{$p$}
\begin{calculs}
p(x) &= \frac{x-1}{2} &&\justif{on recopie l'expression}\\
p(x) &= \frac{x}{2}-\frac{1}{2} &&\justif{on divise chaque terme par $2$}\\
p(x) &= \frac{1}{2}x-\frac{1}{2} &&
\end{calculs}
\explique{$p$ est affine : $\res{a=\frac{1}{2}}$ et $\res{b=-\frac{1}{2}}$.}{}
\rem{Erreur fréquente pour $p$ : oublier de diviser le $-1$ par $2$ et répondre $b=-1$.}
\end{exo}

\begin{exo}[Développer d'abord]
\cas{$f$}
\begin{calculs}
f(x) &= (x+3)^{2}-x^{2} &&\justif{on recopie l'expression}\\
f(x) &= x^{2}+6x+9-x^{2} &&\justif{identité remarquable $(a+b)^{2}=a^{2}+2ab+b^{2}$}\\
f(x) &= 6x+9 &&\justif{les $x^{2}$ s'annulent}
\end{calculs}
\explique{$f$ est affine, avec $a=6$ et $b=9$ ; $a>0$, donc $f$ est \res{croissante sur
$\R$}.}{propriété 2 : le signe de $a$ donne les variations}

\cas{$g$}
\begin{calculs}
g(x) &= (2x-1)(x+2)-2x^{2} &&\justif{on recopie l'expression}\\
g(x) &= 2x^{2}+4x-x-2-2x^{2} &&\justif{double distributivité}\\
g(x) &= 3x-2 &&\justif{on réduit : $4x-x=3x$, les $x^{2}$ s'annulent}
\end{calculs}
\explique{$g$ est affine, avec $a=3$ et $b=-2$ ; $a>0$, donc $g$ est \res{croissante sur
$\R$}.}{}

\cas{$h$}
\begin{calculs}
h(x) &= (x+1)(x-1)-x(x+2) &&\justif{on recopie l'expression}\\
h(x) &= x^{2}-1-\left(x^{2}+2x\right) &&\justif{$(a+b)(a-b)=a^{2}-b^{2}$ ; on garde la
parenthèse}\\
h(x) &= x^{2}-1-x^{2}-2x &&\justif{le signe $-$ change les signes dans la parenthèse}\\
h(x) &= -2x-1 &&\justif{on réduit}
\end{calculs}
\explique{$h$ est affine, avec $a=-2$ et $b=-1$ ; $a<0$, donc $h$ est \res{décroissante sur
$\R$}.}{}
\rem{Sous sa première forme, aucune de ces fonctions n'a l'air affine : il faut développer
et réduire pour le voir. Pour $h$, erreur fréquente : écrire $-x^{2}+2x$ en oubliant que le
$-$ porte sur toute la parenthèse.}
\end{exo}

\begin{exo}[Sens de variation]
\cas{1)}
\explique{$f(x)=-3x+1$ : $a=-3<0$, donc $f$ est \res{décroissante}.}{propriété 2}
\explique{$g(x)=0{,}5x-4$ : $a=0{,}5>0$, donc $g$ est \res{croissante}.}{}
\explique{$h(x)=7$ : $a=0$, donc $h$ est \res{constante}.}{sa droite est horizontale}
\explique{$k(x)=2-\dfrac{x}{5}$, soit $k(x)=-\dfrac{1}{5}x+2$ : $a=-\dfrac15<0$, donc $k$
est \res{décroissante}.}{on remet sous la forme $ax+b$ pour lire $a$}

\cas{2)}
\explique{$-2<3$ et $f$ est décroissante, donc $\res{f(-2)\ge f(3)}$.}{l'ordre est
renversé}
\explique{$-2<3$ et $g$ est croissante, donc $\res{g(-2)\le g(3)}$.}{l'ordre est conservé}
\rem{Contrôle par le calcul (non demandé) : $f(-2)=6+1=7$ et $f(3)=-9+1=-8$ ;
$g(-2)=-1-4=-5$ et $g(3)=1{,}5-4=-2{,}5$.}
\end{exo}

\begin{exo}[Tracer des droites]
\explique{Une droite est tracée dès qu'on connaît deux de ses points : on calcule deux
images.}{méthode 6 du cours}
\begin{lecture}
$f(0)=-1$ et $f(2)=4-1=3$ : points $(0\,;-1)$ et $(2\,;3)$.

\smallskip
$g(0)=3$ et $g(3)=-3+3=0$ : points $(0\,;3)$ et $(3\,;0)$.

\smallskip
$h(0)=1$ et $h(2)=1+1=2$ : points $(0\,;1)$ et $(2\,;2)$.

\smallskip
\gris{pour $h$, on choisit $x$ pair : la moitié de $x$ est entière}

\smallskip
$k(x)=-2$ pour tout $x$ : droite horizontale passant par $(0\,;-2)$.

\smallskip
\gris{Contrôle : $f$ monte de $2$ quand $x$ avance de $1$ ($a=2$) ; $g$ descend de $1$
($a=-1$).}
\end{lecture}
\begin{figlecture}
\begin{tikzpicture}[x=0.4cm,y=0.4cm]
  \repere{-3}{-4}{5}{5}
  \gradx{-2,-1,1,2,3,4} \grady{-3,-2,-1,1,2,3,4}
  \draw[coulCourbe,line width=1pt] (-1.5,-4)--(3,5);
  \draw[coulCourbeB,line width=1pt] (-2,5)--(5,-2);
  \draw[coulCourbeC,line width=1pt] (-3,-0.5)--(5,3.5);
  \draw[black!70,line width=1pt] (-3,-2)--(5,-2);
  \path[coulCourbe] (0,-1) node[croix]{} (2,3) node[croix]{};
  \path[coulCourbeB] (0,3) node[croix]{} (3,0) node[croix]{};
  \path[coulCourbeC] (0,1) node[croix]{} (2,2) node[croix]{};
  \node[coulCourbe,font=\footnotesize,right] at (2.9,4.6) {$f$};
  \node[coulCourbeB,font=\footnotesize,right] at (-1.9,4.6) {$g$};
  \node[coulCourbeC,font=\footnotesize,above] at (4.5,3.3) {$h$};
  \node[black!70,font=\footnotesize,below] at (4.5,-2) {$k$};
\end{tikzpicture}
\end{figlecture}
\end{exo}

\begin{exo}[Deux formules d'abonnement]
\cas{1)}
\explique{$\res{A(x)=30}$.}{le prix de la formule A ne dépend pas du nombre de séances}
\explique{$\res{B(x)=3x+12}$.}{$3$~€ par séance, soit $3x$ pour $x$ séances, plus $12$~€
fixes}

\cas{2)}
\explique{$A$ est \res{constante} ($a=0$ et $b=30$).}{droite horizontale}
\explique{$B$ est \res{affine} ($a=3$ et $b=12$) ; $a>0$, donc $B$ est \res{croissante}.}{plus
on fait de séances, plus on paie}

\cas{3)}
\explique{$A$ : droite horizontale à la hauteur $30$. $B$ : $B(0)=12$ et
$B(10)=30+12=42$, points $(0\,;12)$ et $(10\,;42)$.}{en ordonnée, un carreau vaut $5$~€ :
$30$~€ est au 6\textsuperscript{e} carreau}
\begin{center}
\begin{tikzpicture}[x=0.45cm,y=0.09cm]
  \draw[grille,xstep=1,ystep=5] (0,0) grid (10,45);
  \draw[-{Stealth[length=2mm]},thick] (0,0)--(11,0) node[right,font=\footnotesize]{$x$};
  \draw[-{Stealth[length=2mm]},thick] (0,0)--(0,48) node[right,font=\footnotesize]{prix (€)};
  \foreach \k in {1,...,10} \draw (\k,1)--(\k,-1) node[below,font=\scriptsize,inner sep=1.5pt]{$\k$};
  \foreach \k in {10,20,...,40} \draw (0.15,\k)--(-0.15,\k) node[left,font=\scriptsize,inner sep=1.5pt]{$\k$};
  \node[below left,font=\scriptsize,inner sep=1.5pt] at (0,0) {$0$};
  \draw[coulCourbe,line width=1pt] (0,30)--(10,30) node[above left,font=\footnotesize]{$A$};
  \draw[coulCourbeB,line width=1pt] (0,12)--(10,42) node[left,font=\footnotesize]{$B$};
  \draw[black!50,dashed] (6,30)--(6,0);
  \path (6,30) node[croix]{};
  \morceau{7}{10}
\end{tikzpicture}
\end{center}

\cas{4)}
\explique{Les deux droites se coupent au point d'abscisse $6$. À droite de ce point, la
droite de $B$ est au-dessus de celle de $A$ : la formule B coûte plus cher.}{« moins
cher » : la droite la plus basse}
\explique{La formule A est la moins chère \res{à partir de $7$ séances}.}{$x$ est un nombre
entier ; à $6$ séances, les prix sont égaux (en rouge sur l'axe : de $7$ à $10$ séances)}

\cas{5)}
\begin{calculs}
B(6) &= 3\times 6+12 &&\justif{on remplace $x$ par $6$}\\
B(6) &= 30 &&
\end{calculs}
\explique{$B(6)=A(6)$ : à $6$ séances, les deux formules coûtent \res{$30$~€}, et au-delà B
coûte plus que $30$~€.}{$B$ est croissante : $B(7)=33$, $B(8)=36$\ldots}
\end{exo}

\partie{Fonction affine par lecture graphique}

\begin{exo}[Lire l'expression]
\cas{1)}
\explique{Pour chaque droite : $b$ se lit là où elle coupe l'axe des ordonnées ; pour $a$,
on part de ce point et on avance jusqu'à un autre nœud du quadrillage situé sur la
droite.}{$a=\dfrac{\text{déplacement vertical}}{\text{déplacement horizontal}}$ (méthode 7)}
\begin{lecture}
$(d_1)$ : $b=-2$. De $(0\,;-2)$ à $(1\,;0)$, on avance de $1$ et on monte de $2$ :
$a=\dfrac{2}{1}$, soit $a=2$. Donc $\res{f_1(x)=2x-2}$.

\medskip
$(d_2)$ : $b=2$. De $(0\,;2)$ à $(2\,;1)$, on avance de $2$ et on descend de $1$ :
$a=\dfrac{-1}{2}$. Donc $\res{f_2(x)=-\frac{1}{2}x+2}$.

\smallskip
\gris{en avançant de $1$, on ne tombe pas sur un nœud : on avance de $2$}

\medskip
$(d_3)$ : $b=-1$. De $(0\,;-1)$ à $(1\,;-2)$, on avance de $1$ et on descend de $1$ :
$a=-1$. Donc $\res{f_3(x)=-x-1}$.

\medskip
$(d_4)$ : $b=0$. De $(0\,;0)$ à $(3\,;2)$, on avance de $3$ et on monte de $2$ :
$a=\dfrac{2}{3}$. Donc $\res{f_4(x)=\frac{2}{3}x}$.
\end{lecture}
\begin{figlecture}
\begin{tikzpicture}[x=0.45cm,y=0.45cm]
  \droitesXX
  \draw[coulCourbe,thick,dashed] (0,-2)--(1,-2)--(1,0);
  \draw[coulCourbeB,thick,dashed] (0,2)--(2,2)--(2,1);
  \draw[coulCourbeC,thick,dashed] (0,-1)--(1,-1)--(1,-2);
  \draw[black!70,thick,dashed] (0,0)--(3,0)--(3,2);
\end{tikzpicture}
\end{figlecture}
\rem{Le déplacement vertical est négatif quand on descend : c'est le cas pour $(d_2)$ et
$(d_3)$, et leur coefficient $a$ est négatif.}

\cas{2)}
\explique{\res{$f_4$} est linéaire.}{$b=0$ : sa droite passe par l'origine du repère}
\end{exo}

\begin{exo}[Associer fonctions et droites]
\explique{On commence par le signe de $a$, puis on regarde $b$.}{$a>0$ : la droite monte ;
$a<0$ : elle descend ; $b$ : point où elle coupe l'axe des ordonnées}
\explique{$f$ et $g$ ont $a=2>0$ : leurs droites montent. Ce sont $(D_2)$ et
$(D_4)$.}{$(D_1)$ et $(D_3)$ descendent}
\explique{$f(0)=1$ : \res{$f$ est représentée par $(D_2)$}, qui passe par $(0\,;1)$.}{}
\explique{$g(0)=-3$ : \res{$g$ est représentée par $(D_4)$}, qui passe par $(0\,;-3)$.}{les
deux droites sont parallèles : même $a$}
\explique{$h$ et $k$ ont $a<0$ et $b=1$ : leurs droites descendent et passent toutes deux
par $(0\,;1)$.}{il faut un autre critère : la pente}
\explique{$h$ descend de $1$ quand $x$ avance de $1$ ($a=-1$) : \res{$h$ est représentée par
$(D_1)$}, qui passe par $(1\,;0)$.}{$h(1)=0$}
\explique{$k$ descend de $1$ quand $x$ avance de $3$ ($a=-\frac13$) : \res{$k$ est
représentée par $(D_3)$}, qui passe par $(3\,;0)$.}{$k(3)=-1+1=0$ ; elle descend moins
vite}
\end{exo}

\begin{exo}[Location de vélo]
\cas{1)}
\explique{Pour $2$~h, le point de la droite d'abscisse $2$ a pour ordonnée $10$ : le prix est
de \res{$10$~€}.}{on monte de $x=2$ jusqu'à la droite, puis on lit sur l'axe des
ordonnées (un carreau vaut $2$~€)}

\cas{2)}
\explique{La droite coupe l'axe des ordonnées en $(0\,;4)$ : $\res{b=4}$.}{ordonnée à
l'origine}
\explique{C'est le prix pour une durée nulle : \res{un prix fixe de $4$~€}, payé quelle que
soit la durée.}{frais de location, par exemple}

\cas{3)}
\explique{De $(0\,;4)$ à $(2\,;10)$, on avance de $2$ et on monte de $6$ :
$a=\dfrac{6}{2}$, soit $\res{a=3}$.}{deux nœuds du quadrillage sur la droite}
\explique{Le prix augmente de $3$~€ par heure : \res{$3$~€ est le prix d'une heure de
location}.}{quand $x$ augmente de $1$, $p(x)$ augmente de $a$}

\cas{4)}
\explique{$\res{p(x)=3x+4}$.}{$p(x)=ax+b$ avec $a=3$ et $b=4$}
\begin{calculs}
p(8) &= 3\times 8+4 &&\justif{on remplace $x$ par $8$}\\
p(8) &= 28 &&
\end{calculs}
\explique{Pour $8$~h, le prix est de \res{$28$~€}.}{$8$~h est hors du graphique : on
calcule. Contrôle avec 1) : $p(2)=6+4=10$}
\end{exo}

\partie{Fonction affine passant par deux points}

\begin{exo}[Deux images connues]
\cas{$f$}
\explique{La droite passe par $A(1\,;4)$ et $B(3\,;10)$.}{$f(1)=4$ : point d'abscisse $1$ et
d'ordonnée $4$}
\begin{calculs}
a &= \frac{y_{B}-y_{A}}{x_{B}-x_{A}} &&\justif{propriété 3}\\
a &= \frac{10-4}{3-1} &&\justif{on remplace}\\
a &= 3 &&\justif{$6\div 2$}
\end{calculs}
\begin{calculs}
f(1) &= 3\times 1+b &&\justif{on remplace $x$ par $1$ dans $f(x)=3x+b$}\\
4 &= 3+b &&\justif{car $f(1)=4$}\\
b &= 1 &&\justif{on enlève $3$ des deux côtés}
\end{calculs}
\explique{$\res{f(x)=3x+1}$.}{contrôle avec $B$ : $f(3)=9+1=10$}

\cas{$g$}
\explique{La droite passe par $A(-2\,;7)$ et $B(2\,;-1)$.}{}
\begin{calculs}
a &= \frac{-1-7}{2-(-2)} &&\justif{propriété 3}\\
a &= \frac{-8}{4} &&\justif{$2-(-2)=2+2$}\\
a &= -2 &&
\end{calculs}
\begin{calculs}
g(2) &= -2\times 2+b &&\justif{on remplace $x$ par $2$}\\
-1 &= -4+b &&\justif{car $g(2)=-1$}\\
b &= 3 &&\justif{on ajoute $4$ des deux côtés}
\end{calculs}
\explique{$\res{g(x)=-2x+3}$.}{contrôle avec $A$ : $g(-2)=4+3=7$}

\cas{$h$}
\explique{$h(0)=5$ : c'est directement l'ordonnée à l'origine, $b=5$.}{$h(0)=a\times 0+b$,
soit $b$}
\begin{calculs}
a &= \frac{3-5}{4-0} &&\justif{propriété 3, avec $(0\,;5)$ et $(4\,;3)$}\\
a &= \frac{-2}{4} &&\\
a &= -\frac{1}{2} &&\justif{on simplifie par $2$}
\end{calculs}
\explique{$\res{h(x)=-\frac{1}{2}x+5}$.}{contrôle : $h(4)=-2+5=3$}

\cas{$k$}
\explique{La droite passe par $A(-3\,;-4)$ et $B(6\,;2)$.}{}
\begin{calculs}
a &= \frac{2-(-4)}{6-(-3)} &&\justif{propriété 3}\\
a &= \frac{6}{9} &&\justif{$2+4=6$ et $6+3=9$}\\
a &= \frac{2}{3} &&\justif{on simplifie par $3$}
\end{calculs}
\begin{calculs}
k(6) &= \frac{2}{3}\times 6+b &&\justif{on prend $B$ : $\frac23\times 6$ est entier}\\
2 &= 4+b &&\justif{car $k(6)=2$}\\
b &= -2 &&\justif{on enlève $4$ des deux côtés}
\end{calculs}
\explique{$\res{k(x)=\frac{2}{3}x-2}$.}{contrôle avec $A$ : $k(-3)=-2-2=-4$}
\rem{Erreur fréquente : les soustractions de nombres négatifs. $2-(-4)=2+4$, et non
$2-4$.}
\end{exo}

\begin{exo}[Droite passant par deux points]
\cas{1)}
\begin{calculs}
a &= \frac{3-1}{6-2} &&\justif{propriété 3, avec $A(2\,;1)$ et $B(6\,;3)$}\\
a &= \frac{2}{4} &&\\
a &= \frac{1}{2} &&
\end{calculs}
\begin{calculs}
f(2) &= \frac{1}{2}\times 2+b &&\justif{on remplace $x$ par $2$}\\
1 &= 1+b &&\justif{car $f(2)=1$}\\
b &= 0 &&
\end{calculs}
\explique{$\res{f(x)=\frac{1}{2}x}$ : $b=0$, la fonction est \res{linéaire}.}{sa droite passe
par l'origine du repère}

\cas{2)}
\begin{calculs}
f(10) &= \frac{1}{2}\times 10 &&\justif{image de l'abscisse de $E$}\\
f(10) &= 5 &&
\end{calculs}
\explique{$f(10)=5$, l'ordonnée de $E$ : \res{oui, $E$ est sur la droite $(AB)$}.}{on
compare l'image de l'abscisse à l'ordonnée du point}

\cas{3)}
\begin{calculs}
a &= \frac{-4-5}{2-(-1)} &&\justif{propriété 3, avec $C(-1\,;5)$ et $D(2\,;-4)$}\\
a &= \frac{-9}{3} &&\justif{$2-(-1)=2+1$}\\
a &= -3 &&
\end{calculs}
\begin{calculs}
g(-1) &= -3\times(-1)+b &&\justif{on remplace $x$ par $-1$}\\
5 &= 3+b &&\justif{car $g(-1)=5$}\\
b &= 2 &&
\end{calculs}
\explique{$\res{g(x)=-3x+2}$.}{contrôle avec $D$ : $g(2)=-6+2=-4$}

\cas{4)}
\begin{calculs}
g(3) &= -3\times 3+2 &&\justif{image de l'abscisse de $F$}\\
g(3) &= -7 &&
\end{calculs}
\explique{$g(3)=-7$, et $-7\neq -6$ : \res{non, $F$ n'est pas sur la droite $(CD)$}.}{le point
de la droite d'abscisse $3$ est $(3\,;-7)$, un peu plus bas que $F$}
\end{exo}

\begin{exo}[Température]
\cas{1)}
\explique{$\res{T(8)=6}$ et $\res{T(14)=18}$.}{$6$~°C à 8~h, $18$~°C à 14~h}
\begin{calculs}
a &= \frac{18-6}{14-8} &&\justif{propriété 3, avec $(8\,;6)$ et $(14\,;18)$}\\
a &= \frac{12}{6} &&\\
a &= 2 &&\justif{la température monte de $2$~°C par heure}
\end{calculs}
\begin{calculs}
T(8) &= 2\times 8+b &&\justif{on remplace $h$ par $8$}\\
6 &= 16+b &&\justif{car $T(8)=6$}\\
b &= -10 &&\justif{on enlève $16$ des deux côtés}
\end{calculs}
\explique{$\res{T(h)=2h-10}$.}{contrôle : $T(14)=28-10=18$}

\cas{2)}
\begin{calculs}
T(11) &= 2\times 11-10 &&\justif{on remplace $h$ par $11$}\\
T(11) &= 12 &&
\end{calculs}
\explique{À 11~h, il fait \res{$12$~°C}.}{11~h est au milieu de 8~h et 14~h, et $12$ au
milieu de $6$ et $18$}

\cas{3)}
\explique{On cherche $h$ tel que $T(h)=14$ : c'est une équation.}{on connaît l'image, on
cherche l'antécédent}
\begin{calculs}
2h-10 &= 14 &&\justif{on recopie l'équation}\\
2h &= 24 &&\justif{on ajoute $10$}\\
h &= 12 &&\justif{on divise par $2$}
\end{calculs}
\explique{Il fait $14$~°C \res{à 12~h}.}{contrôle : $T(12)=24-10=14$}
\rem{$b=-10$ n'a pas de sens concret ici (il ne fait pas $-10$~°C à minuit) : le modèle
affine ne vaut que sur $\intff{8}{14}$.}
\end{exo}

\begin{exo}[Facture d'électricité]
\cas{1)}
\explique{La droite passe par $(300\,;72)$ et $(500\,;102)$.}{consommation en abscisse, prix
en ordonnée}
\begin{calculs}
a &= \frac{102-72}{500-300} &&\justif{propriété 3}\\
a &= \frac{30}{200} &&\\
a &= 0{,}15 &&\justif{$30\div 200$}
\end{calculs}
\begin{calculs}
f(300) &= 0{,}15\times 300+b &&\justif{on remplace $x$ par $300$}\\
72 &= 45+b &&\justif{car $f(300)=72$}\\
b &= 27 &&\justif{on enlève $45$ des deux côtés}
\end{calculs}
\explique{$\res{f(x)=0{,}15x+27}$.}{contrôle : $f(500)=75+27=102$}

\cas{2)}
\explique{$a=0{,}15$ : \res{le prix d'un kWh}, $0{,}15$~€.}{chaque kWh de plus ajoute
$0{,}15$~€ à la facture}
\explique{$b=27$ : \res{l'abonnement}, $27$~€, payé même sans consommer.}{$f(0)=27$}

\cas{3)}
\begin{calculs}
f(800) &= 0{,}15\times 800+27 &&\justif{on remplace $x$ par $800$}\\
f(800) &= 120+27 &&\\
f(800) &= 147 &&
\end{calculs}
\explique{Pour $800$~kWh, on paie \res{$147$~€}.}{}
\end{exo}

\partie{Signe d'une fonction par lecture graphique}

\begin{exo}[Signe d'une fonction]
\cas{1)}
\begin{lecture}
Résoudre $f(x)=0$, c'est chercher les abscisses des points de $\Cf$ d'ordonnée $0$ : les
points où la courbe coupe l'axe des abscisses.

\smallskip
$\Cf$ coupe l'axe des abscisses en $x=-3$, en $x=1$ et en $x=4$ :
$\res{S=\ens{-3\,;1\,;4}}$.

\smallskip
\gris{les croix rouges sur l'axe ; on écrit les solutions entre accolades : c'est un
ensemble de trois nombres, pas un intervalle}
\end{lecture}
\begin{figlecture}
\begin{tikzpicture}[x=0.42cm,y=0.42cm]
  \courbeXXVII
  \solx{-3}\solx{1}\solx{4}
\end{tikzpicture}
\end{figlecture}

\cas{2)}
\explique{Première ligne : les bornes $-4$ et $5$, et les valeurs où $f$ s'annule, $-3$, $1$
et $4$, rangées dans l'ordre.}{définition 5 du cours}
\explique{Seconde ligne : un $0$ sous chaque valeur qui annule $f$ ; entre deux, $+$ si la
courbe est au-dessus de l'axe, $-$ si elle est au-dessous.}{de gauche à droite :
au-dessus, au-dessous, au-dessus, au-dessous}
\begin{center}
\begin{tikzpicture}[scale=0.85]
  \tkzTabInit[lgt=1.4,espcl=1.6]{$x$/1,$f(x)$/1}{$-4$,$-3$,$1$,$4$,$5$}
  \tkzTabLine{,+,z,-,z,+,z,-,}
\end{tikzpicture}
\end{center}

\cas{3)}
\begin{lecture}
$f(x)>0$ : la courbe est \textbf{au-dessus} de l'axe des abscisses. C'est le cas de
$x=-4$ à $x=-3$, puis de $x=1$ à $x=4$ (en rouge et en bleu sur l'axe).

\smallskip
$\res{S=\intfo{-4}{-3}\cup\intoo{1}{4}}$.

\smallskip
\gris{Ce sont les signes $+$ du tableau de 2). Inégalité stricte : $-3$, $1$ et $4$ sont
exclus, car $f$ y vaut $0$. En revanche $-4$ est compris : $f(-4)=2$, et $2>0$.}

\smallskip
\gris{Le symbole $\cup$ (« union ») réunit les deux intervalles.}
\end{lecture}
\begin{figlecture}
\begin{tikzpicture}[x=0.42cm,y=0.42cm]
  \courbeXXVII
  \morceau{-4}{-3}\crochet{-4}{[}\crochet{-3}{[}
  \morceau[coulIntB]{1}{4}\crochet[coulIntB]{1}{]}\crochet[coulIntB]{4}{[}
\end{tikzpicture}
\end{figlecture}

\cas{4)}
\explique{Première ligne : $-4$ ; $-1$ ; $2$ ; $5$.}{bornes, puis changements de sens : le
creux $(-1\,;-2)$ et le sommet $(2\,;2)$}
\explique{La courbe descend de $2$ à $-2$, monte jusqu'à $2$, puis redescend jusqu'à
$-2$.}{images au bout des flèches lues sur la courbe}
\begin{center}
\begin{tikzpicture}[scale=0.85]
  \tkzTabInit[lgt=1.4,espcl=2,deltacl=0.6]{$x$/1,$f(x)$/2}{$-4$,$-1$,$2$,$5$}
  \tkzTabVar{+/$2$,-/$-2$,+/$2$,-/$-2$}
\end{tikzpicture}
\end{center}
\rem{Ne pas confondre les deux tableaux. Le tableau de signes dit si la courbe est
au-dessus ou au-dessous de l'axe des abscisses (on regarde où elle le coupe) ; le tableau de
variations dit si elle monte ou descend (on regarde ses sommets et ses creux).}
\end{exo}

\begin{exo}[Signe de fonctions affines]
\explique{Pour chaque droite de l'exercice 20, on lit où elle coupe l'axe des abscisses,
puis si elle est au-dessus ou au-dessous de l'axe de chaque côté.}{méthode 9 du cours ; une
fonction affine est définie sur $\R$ : première ligne de $-\infty$ à $+\infty$}

\cas{$f_1$}
\explique{$(d_1)$ coupe l'axe en $x=1$ ; au-dessous avant, au-dessus après.}{la droite
monte : $a=2>0$}
\begin{center}\signeaffine{$f_1(x)$}{$1$}{-}{+}\end{center}

\cas{$f_2$}
\explique{$(d_2)$ coupe l'axe en $x=4$ ; au-dessus avant, au-dessous après.}{la droite
descend : $a=-\frac12<0$}
\begin{center}\signeaffine{$f_2(x)$}{$4$}{+}{-}\end{center}

\cas{$f_3$}
\explique{$(d_3)$ coupe l'axe en $x=-1$ ; au-dessus avant, au-dessous après.}{la droite
descend : $a=-1<0$}
\begin{center}\signeaffine{$f_3(x)$}{$-1$}{+}{-}\end{center}

\cas{$f_4$}
\explique{$(d_4)$ passe par l'origine : elle coupe l'axe en $x=0$ ; au-dessous avant,
au-dessus après.}{la droite monte : $a=\frac23>0$}
\begin{center}\signeaffine{$f_4(x)$}{$0$}{-}{+}\end{center}
\rem{Contrôle par le calcul avec les expressions de l'exercice 20 : $f_1(1)=2-2=0$ et
$f_2(4)=-2+2=0$.}
\end{exo}

\begin{exo}[Des variations au signe]
\explique{Sur $\intff{-3}{1}$, $f$ est croissante, de $-2$ à $4$ : elle passe par $0$ une
seule fois, en $-1$.}{elle part d'une valeur négative et arrive à une valeur positive}
\explique{Donc $f(x)<0$ de $-3$ à $-1$, et $f(x)>0$ de $-1$ à $1$.}{avant $-1$, $f(x)$ est
plus petit que $f(-1)=0$ ; après, il est plus grand}
\explique{Sur $\intff{1}{5}$, $f$ est décroissante, de $4$ à $-1$ : elle passe par $0$ une
seule fois, en $3$.}{}
\explique{Donc $f(x)>0$ de $1$ à $3$, et $f(x)<0$ de $3$ à $5$.}{décroissante : avant $3$,
$f(x)$ est plus grand que $f(3)=0$}
\begin{center}
\begin{tikzpicture}[scale=0.85]
  \tkzTabInit[lgt=1.4,espcl=2]{$x$/1,$f(x)$/1}{$-3$,$-1$,$3$,$5$}
  \tkzTabLine{,-,z,+,z,-,}
\end{tikzpicture}
\end{center}
\rem{La valeur $1$, où $f$ change de sens, n'apparaît pas dans le tableau de signes :
$f(1)=4$ n'est pas nul. On n'y écrit que les bornes et les valeurs qui annulent $f$.}
\end{exo}

\partie{Propriétés des inégalités, encadrer une quantité}

\begin{exo}[Règles des inégalités]
\explique{a) $\res{a+4<b+4}$.}{ajouter un nombre conserve le sens (propriété 4)}
\explique{b) $\res{a-7<b-7}$.}{soustraire un nombre conserve le sens}
\explique{c) $\res{3a<3b}$.}{multiplier par $3>0$ conserve le sens}
\explique{d) $\res{-5a>-5b}$.}{multiplier par $-5<0$ change le sens}
\explique{e) $\res{\dfrac{a}{2}<\dfrac{b}{2}}$.}{diviser par $2>0$ conserve le sens}
\explique{f) $\res{-a>-b}$.}{$-a$, c'est $(-1)\times a$ : on multiplie par $-1<0$, le sens
change}
\rem{Contrôle avec des nombres : $a=1$ et $b=2$. En d), $-5>-10$ ; en f), $-1>-2$. Sur la
droite graduée, les opposés sont rangés dans l'autre sens.}
\end{exo}

\begin{exo}[Encadrer]
\explique{On part de l'encadrement de $x$ et on construit l'expression étape par étape, en
appliquant la même opération aux trois membres.}{méthode 10 du cours}

\cas{$A$}
\begin{calculs}
\renewcommand{\arraystretch}{1.5}
\begin{array}{@{}r@{\;}c@{\;}c@{\;}c@{\;}l@{\qquad}l@{}}
1 & \le & x & \le & 4 & \justif{on recopie l'encadrement}\\
2 & \le & 2x & \le & 8 & \justif{on multiplie par $2>0$ : même sens}\\
5 & \le & 2x+3 & \le & 11 & \justif{on ajoute $3$}
\end{array}
\end{calculs}
\explique{$\res{5\le A\le 11}$.}{}

\cas{$B$}
\begin{calculs}
\renewcommand{\arraystretch}{1.5}
\begin{array}{@{}r@{\;}c@{\;}c@{\;}c@{\;}l@{\qquad}l@{}}
1 & \le & x & \le & 4 & \justif{on recopie l'encadrement}\\
-3 & \ge & -3x & \ge & -12 & \justif{on multiplie par $-3<0$ : le sens change}\\
-2 & \ge & -3x+1 & \ge & -11 & \justif{on ajoute $1$}
\end{array}
\end{calculs}
\explique{$\res{-11\le B\le -2}$.}{on réécrit l'encadrement dans l'ordre croissant}

\cas{$C$}
\begin{calculs}
\renewcommand{\arraystretch}{1.5}
\begin{array}{@{}r@{\;}c@{\;}c@{\;}c@{\;}l@{\qquad}l@{}}
1 & \le & x & \le & 4 & \justif{on recopie l'encadrement}\\
-2 & \ge & -2x & \ge & -8 & \justif{on multiplie par $-2<0$ : le sens change}\\
3 & \ge & 5-2x & \ge & -3 & \justif{on ajoute $5$}
\end{array}
\end{calculs}
\explique{$\res{-3\le C\le 3}$.}{$5-2x$, c'est $-2x+5$}

\cas{$D$}
\begin{calculs}
\renewcommand{\arraystretch}{1.9}
\begin{array}{@{}r@{\;}c@{\;}c@{\;}c@{\;}l@{\qquad}l@{}}
1 & \le & x & \le & 4 & \justif{on recopie l'encadrement}\\
\dfrac{1}{2} & \le & \dfrac{1}{2}x & \le & 2 & \justif{on multiplie par $\frac12>0$ : même
sens}\\
-\dfrac{1}{2} & \le & \dfrac{1}{2}x-1 & \le & 1 & \justif{on enlève $1$ : $\frac12-1=-\frac12$}
\end{array}
\end{calculs}
\explique{$\res{-\dfrac{1}{2}\le D\le 1}$.}{}
\rem{Après une multiplication par un nombre négatif, le plus grand nombre passe à gauche :
on réécrit l'encadrement dans l'ordre croissant. Erreur fréquente pour $B$ : garder
$\le$ et écrire $-3\le -3x\le -12$, ce qui est impossible ($-3>-12$).}
\end{exo}

\begin{exo}[Mesures imprécises]
\cas{1)}
\explique{Périmètre :}{on part de l'encadrement de la longueur $x$}
\begin{calculs}
\renewcommand{\arraystretch}{1.5}
\begin{array}{@{}r@{\;}c@{\;}c@{\;}c@{\;}l@{\qquad}l@{}}
10 & \le & x & \le & 12 & \justif{on recopie l'encadrement}\\
20 & \le & 2x & \le & 24 & \justif{on multiplie par $2>0$ : même sens}\\
32 & \le & 2x+12 & \le & 36 & \justif{on ajoute $12$}
\end{array}
\end{calculs}
\explique{Le périmètre est compris entre $32$~m et $36$~m : $\res{32\le P\le 36}$.}{$P=2x+2\times
6$}
\explique{Aire :}{}
\begin{calculs}
\renewcommand{\arraystretch}{1.5}
\begin{array}{@{}r@{\;}c@{\;}c@{\;}c@{\;}l@{\qquad}l@{}}
10 & \le & x & \le & 12 & \justif{on recopie l'encadrement}\\
60 & \le & 6x & \le & 72 & \justif{on multiplie par $6>0$ : même sens}
\end{array}
\end{calculs}
\explique{L'aire est comprise entre $60$~m² et $72$~m² : $\res{60\le \mathcal{A}\le
72}$.}{longueur fois largeur}

\cas{2)}
\begin{calculs}
\renewcommand{\arraystretch}{1.5}
\begin{array}{@{}r@{\;}c@{\;}c@{\;}c@{\;}l@{\qquad}l@{}}
15 & \le & C & \le & 25 & \justif{on recopie l'encadrement}\\
27 & \le & 1{,}8C & \le & 45 & \justif{on multiplie par $1{,}8>0$ : même sens}\\
59 & \le & 1{,}8C+32 & \le & 77 & \justif{on ajoute $32$}
\end{array}
\end{calculs}
\explique{Au printemps, la température est comprise entre $59$ et $77$ degrés Fahrenheit :
$\res{59\le F\le 77}$.}{$1{,}8\times 15=15+12$ et $1{,}8\times 25=25+20$}
\end{exo}

\partie{Résoudre une inéquation du premier degré}

\begin{exo}[Est-ce une solution ?]
\explique{On remplace $x$ par $3$ dans le membre de gauche (et dans le membre de droite s'il
contient $x$), puis on regarde si l'inégalité obtenue est vraie.}{comme pour vérifier une
solution d'équation}
\explique{a) $2\times 3-1=5$, et $5>4$ est vrai : \res{$3$ est solution}.}{}
\explique{b) $-3+5=2$, et $2\ge 2$ est vrai : \res{$3$ est solution}.}{« $\ge$ » :
l'égalité convient}
\explique{c) À gauche : $4\times 3-9=3$ ; à droite : $3$. Or $3<3$ est faux : \res{$3$ n'est
pas solution}.}{inégalité stricte : l'égalité ne convient pas}
\explique{d) $5-2\times 3=-1$, et $-1\le -1$ est vrai : \res{$3$ est solution}.}{« $\le$ » :
l'égalité convient}
\rem{Avec $\le$ ou $\ge$, un nombre qui donne l'égalité est solution (b et d) ; avec $<$ ou
$>$, il ne l'est pas (c).}
\end{exo}

\begin{exo}[Résoudre]
\explique{Comme pour une équation : les $x$ d'un côté, les nombres de l'autre, puis on
divise par le coefficient de $x$, en changeant le sens s'il est négatif.}{méthode 11 du
cours}

\cas{a)}
\begin{calculs}
x+6 &> 2 &&\justif{on recopie l'inéquation}\\
x &> -4 &&\justif{on enlève $6$ : $2-6=-4$}
\end{calculs}
\explique{$\res{S=\intoo{-4}{+\infty}}$.}{inégalité stricte : $-4$ exclu}

\cas{b)}
\begin{calculs}
-3x &< 15 &&\justif{on recopie l'inéquation}\\
x &> -5 &&\justif{on divise par $-3<0$ : le sens change}
\end{calculs}
\explique{$\res{S=\intoo{-5}{+\infty}}$.}{contrôle avec $x=0$ : $0<15$ est vrai, et $0>-5$}

\cas{c)}
\begin{calculs}
2x-5 &\ge 7 &&\justif{on recopie l'inéquation}\\
2x &\ge 12 &&\justif{on ajoute $5$}\\
x &\ge 6 &&\justif{on divise par $2>0$ : même sens}
\end{calculs}
\explique{$\res{S=\intfo{6}{+\infty}}$.}{« $\ge$ » : $6$ compris}

\cas{d)}
\begin{calculs}
9-4x &> 1 &&\justif{on recopie l'inéquation}\\
-4x &> -8 &&\justif{on enlève $9$}\\
x &< 2 &&\justif{on divise par $-4<0$ : le sens change}
\end{calculs}
\explique{$\res{S=\intoo{-\infty}{2}}$.}{contrôle avec $x=0$ : $9>1$ est vrai, et $0<2$}

\cas{e)}
\begin{calculs}
3x+2 &\le 5x-6 &&\justif{on recopie l'inéquation}\\
3x-5x &\le -6-2 &&\justif{les $x$ à gauche, les nombres à droite}\\
-2x &\le -8 &&\justif{on réduit}\\
x &\ge 4 &&\justif{on divise par $-2<0$ : le sens change}
\end{calculs}
\explique{$\res{S=\intfo{4}{+\infty}}$.}{contrôle avec $x=5$ : $17\le 19$ est vrai}

\cas{f)}
\begin{calculs}
-x+4 &\ge 3x-8 &&\justif{on recopie l'inéquation}\\
-x-3x &\ge -8-4 &&\justif{les $x$ à gauche, les nombres à droite}\\
-4x &\ge -12 &&\justif{on réduit}\\
x &\le 3 &&\justif{on divise par $-4<0$ : le sens change}
\end{calculs}
\explique{$\res{S=\intof{-\infty}{3}}$.}{contrôle avec $x=0$ : $4\ge -8$ est vrai}
\rem{On ne change le sens qu'au moment de diviser (ou de multiplier) par un nombre
négatif ; ajouter ou enlever un nombre ne le change jamais.}
\end{exo}

\begin{exo}[Parenthèses et fractions]
\cas{a)}
\begin{calculs}
2(x-3) &< x+1 &&\justif{on recopie l'inéquation}\\
2x-6 &< x+1 &&\justif{on développe}\\
2x-x &< 1+6 &&\justif{les $x$ à gauche, les nombres à droite}\\
x &< 7 &&\justif{on réduit}
\end{calculs}
\explique{$\res{S=\intoo{-\infty}{7}}$.}{contrôle avec $x=0$ : $-6<1$ est vrai}

\cas{b)}
\begin{calculs}
3(2-x) &\ge 2(x+8) &&\justif{on recopie l'inéquation}\\
6-3x &\ge 2x+16 &&\justif{on développe les deux membres}\\
-3x-2x &\ge 16-6 &&\justif{les $x$ à gauche, les nombres à droite}\\
-5x &\ge 10 &&\justif{on réduit}\\
x &\le -2 &&\justif{on divise par $-5<0$ : le sens change}
\end{calculs}
\explique{$\res{S=\intof{-\infty}{-2}}$.}{contrôle avec $x=-3$ : $15\ge 10$ est vrai}

\cas{c)}
\begin{calculs}
\frac{x}{2}+1 &> 4 &&\justif{on recopie l'inéquation}\\
\frac{x}{2} &> 3 &&\justif{on enlève $1$}\\
x &> 6 &&\justif{on multiplie par $2>0$ : même sens}
\end{calculs}
\explique{$\res{S=\intoo{6}{+\infty}}$.}{}

\cas{d)}
\begin{calculs}
\frac{2x-1}{3} &\le x &&\justif{on recopie l'inéquation}\\
2x-1 &\le 3x &&\justif{on multiplie les deux membres par $3>0$ : même sens}\\
2x-3x &\le 1 &&\justif{les $x$ à gauche, les nombres à droite}\\
-x &\le 1 &&\justif{on réduit}\\
x &\ge -1 &&\justif{on multiplie par $-1<0$ : le sens change}
\end{calculs}
\explique{$\res{S=\intfo{-1}{+\infty}}$.}{contrôle avec $x=2$ : $1\le 2$ est vrai}
\rem{En d), on multiplie \textbf{tout} le membre de droite par $3$ : $3\times x=3x$. On se
débarrasse ainsi de la fraction avant de regrouper les $x$.}
\end{exo}

\begin{exo}[Corriger les erreurs]
\cas{1)}
\explique{La première étape est juste : on enlève $3$, et $-2x>6$.}{$9-3=6$}
\explique{L'erreur : on a divisé par $-2$ sans changer le sens.}{$-2<0$}
\begin{calculs}
-2x+3 &> 9 &&\justif{on recopie l'inéquation}\\
-2x &> 6 &&\justif{on enlève $3$}\\
x &< -3 &&\justif{on divise par $-2<0$ : le sens change}
\end{calculs}
\explique{$\res{S=\intoo{-\infty}{-3}}$.}{contrôle avec $x=-4$ : $11>9$ est vrai}

\cas{2)}
\explique{L'erreur : on a changé le sens en divisant par $5$, qui est positif.}{le résultat
$-2$ est négatif, mais c'est le signe du diviseur qui compte}
\begin{calculs}
5x &\le -10 &&\justif{on recopie l'inéquation}\\
x &\le -2 &&\justif{on divise par $5>0$ : même sens}
\end{calculs}
\explique{$\res{S=\intof{-\infty}{-2}}$.}{contrôle avec $x=-3$ : $-15\le -10$ est vrai}
\rem{Le sens ne change que si l'on multiplie ou divise par un nombre négatif ; le signe des
autres nombres de l'inéquation ne compte pas.}
\end{exo}

\partie{Signe d'une fonction affine par le calcul}

\begin{exo}[Tableaux de signes]
\explique{Pour chaque expression : on résout $ax+b=0$ pour trouver la valeur qui l'annule,
puis le signe de $a$ donne l'ordre des signes.}{propriété 5 : $a>0$, $-$ puis $+$ (la droite
monte) ; $a<0$, $+$ puis $-$ (elle descend)}

\cas{$A$}
\begin{calculs}
2x-8 &= 0 &&\justif{valeur qui annule $A(x)$}\\
2x &= 8 &&\justif{on ajoute $8$}\\
x &= 4 &&\justif{on divise par $2$}
\end{calculs}
\explique{$a=2>0$ : $-$ avant $4$, $+$ après.}{contrôle : $A(0)=-8$, négatif, et $0$ est
avant $4$}
\begin{center}\signeaffine{$2x-8$}{$4$}{-}{+}\end{center}

\cas{$B$}
\begin{calculs}
-3x+9 &= 0 &&\justif{valeur qui annule $B(x)$}\\
-3x &= -9 &&\justif{on enlève $9$}\\
x &= 3 &&\justif{on divise par $-3$}
\end{calculs}
\explique{$a=-3<0$ : $+$ avant $3$, $-$ après.}{contrôle : $B(0)=9$, positif}
\begin{center}\signeaffine{$-3x+9$}{$3$}{+}{-}\end{center}

\cas{$C$}
\begin{calculs}
5x+15 &= 0 &&\justif{valeur qui annule $C(x)$}\\
5x &= -15 &&\justif{on enlève $15$}\\
x &= -3 &&\justif{on divise par $5$}
\end{calculs}
\explique{$a=5>0$ : $-$ avant $-3$, $+$ après.}{contrôle : $C(0)=15$, positif, et $0$ est
après $-3$}
\begin{center}\signeaffine{$5x+15$}{$-3$}{-}{+}\end{center}

\cas{$D$}
\begin{calculs}
-\frac{1}{2}x+4 &= 0 &&\justif{valeur qui annule $D(x)$}\\
-\frac{1}{2}x &= -4 &&\justif{on enlève $4$}\\
x &= 8 &&\justif{on multiplie par $-2$ : $-4\times(-2)=8$}
\end{calculs}
\explique{$a=-\frac{1}{2}<0$ : $+$ avant $8$, $-$ après.}{contrôle : $D(0)=4$, positif}
\begin{center}\signeaffine{$-\frac{1}{2}x+4$}{$8$}{+}{-}\end{center}

\cas{$E$}
\explique{$7x=0$ pour $x=0$ ; $a=7>0$ : $-$ avant $0$, $+$ après.}{fonction linéaire : sa
droite passe par l'origine ; c'est le signe de $x$, multiplié par $7>0$}
\begin{center}\signeaffine{$7x$}{$0$}{-}{+}\end{center}

\cas{$F$}
\begin{calculs}
-6x+4 &= 0 &&\justif{valeur qui annule $F(x)$}\\
-6x &= -4 &&\justif{on enlève $4$}\\
x &= \frac{-4}{-6} &&\justif{on divise par $-6$}\\
x &= \frac{2}{3} &&\justif{on simplifie par $-2$}
\end{calculs}
\explique{$a=-6<0$ : $+$ avant $\frac{2}{3}$, $-$ après.}{contrôle : $F(0)=4$, positif}
\begin{center}\signeaffine{$-6x+4$}{$\frac{2}{3}$}{+}{-}\end{center}
\rem{Erreur fréquente : prendre le signe de $b$, ou du nombre trouvé, au lieu du signe de
$a$. Pour contrôler, on calcule une image facile, souvent en $0$.}
\end{exo}

\begin{exo}[Variations, droite et signe]
\cas{1)}
\explique{$f(x)=3x-6$ : $a=3>0$, donc $f$ est \res{croissante sur $\R$}.}{propriété 2}
\explique{$g(x)=8-2x$, soit $g(x)=-2x+8$ : $a=-2<0$, donc $g$ est \res{décroissante sur
$\R$}.}{le terme en $x$ d'abord, pour lire $a$}

\cas{2)}
\begin{lecture}
$f(1)=3-6=-3$ et $f(3)=9-6=3$ : points $(1\,;-3)$ et $(3\,;3)$.

\smallskip
$g(1)=8-2=6$ et $g(4)=8-8=0$ : points $(1\,;6)$ et $(4\,;0)$.

\smallskip
\gris{deux points par droite, à coordonnées entières, puis on trace à la règle}

\medskip
\textbf{\color{coulExo}3) vérification sur le graphique} : la droite de $f$ coupe l'axe des
abscisses en $2$ ; elle est au-dessous avant $2$, au-dessus après. La droite de $g$ le coupe
en $4$ ; elle est au-dessus avant $4$, au-dessous après.

\smallskip
\gris{les croix sur l'axe : les valeurs qui annulent $f(x)$ et $g(x)$}
\end{lecture}
\begin{figlecture}
\begin{tikzpicture}[x=0.36cm,y=0.36cm]
  \repere{-1}{-4}{6}{7}
  \gradx{1,2,3,4,5} \grady{-3,-2,-1,1,2,3,4,5,6}
  \draw[coulCourbe,line width=1pt] (0.67,-4)--(4.33,7);
  \draw[coulCourbeB,line width=1pt] (0.5,7)--(6,-4);
  \path[coulCourbe] (1,-3) node[croix]{} (3,3) node[croix]{} (2,0) node[croix]{};
  \path[coulCourbeB] (1,6) node[croix]{} (4,0) node[croix]{};
  \node[coulCourbe,font=\footnotesize,right] at (4.2,6.4) {$f$};
  \node[coulCourbeB,font=\footnotesize,left] at (0.6,6.4) {$g$};
\end{tikzpicture}
\end{figlecture}

\cas{3)}
\begin{calculs}
3x-6 &= 0 &&\justif{valeur qui annule $f(x)$}\\
3x &= 6 &&\justif{on ajoute $6$}\\
x &= 2 &&\justif{on divise par $3$}
\end{calculs}
\explique{$a=3>0$ : $-$ avant $2$, $+$ après.}{la droite de $f$ monte}
\begin{center}\signeaffine{$f(x)$}{$2$}{-}{+}\end{center}
\begin{calculs}
8-2x &= 0 &&\justif{valeur qui annule $g(x)$}\\
-2x &= -8 &&\justif{on enlève $8$}\\
x &= 4 &&\justif{on divise par $-2$}
\end{calculs}
\explique{$a=-2<0$ : $+$ avant $4$, $-$ après.}{la droite de $g$ descend}
\begin{center}\signeaffine{$g(x)$}{$4$}{+}{-}\end{center}
\explique{Le graphique confirme les deux tableaux (voir à côté de 2)).}{le calcul donne la
valeur exacte, le graphique permet de contrôler}
\end{exo}

\begin{exo}[Associer]
\explique{$f(3)=12-12=0$ et $g(3)=6-6=0$ : les deux fonctions s'annulent en $3$.}{la
première ligne ne permet pas de les distinguer : c'est le signe de $a$ qui décide}
\explique{$f(x)=4x-12$ : $a=4>0$, donc $-$ avant $3$, $+$ après : \res{$f$ correspond au
tableau 2}.}{la droite de $f$ monte}
\explique{$g(x)=6-2x$, soit $g(x)=-2x+6$ : $a=-2<0$, donc $+$ avant $3$, $-$ après : \res{$g$
correspond au tableau 1}.}{la droite de $g$ descend}
\rem{Contrôle avec $x=0$ : $f(0)=-12$, négatif, et $g(0)=6$, positif ; or $0$ est avant
$3$.}
\end{exo}

\begin{exo}[Bénéfice]
\cas{1)}
\begin{calculs}
B(0) &= 12\times 0-480 &&\justif{on remplace $x$ par $0$}\\
B(0) &= -480 &&
\end{calculs}
\explique{$\res{B(0)=-480}$ : sans aucune vente, l'artisane perd $480$~€.}{ce sont ses frais
fixes du mois (matériel, local\ldots) ; un bénéfice négatif est une perte}

\cas{2)}
\begin{calculs}
12x-480 &= 0 &&\justif{valeur qui annule $B(x)$}\\
12x &= 480 &&\justif{on ajoute $480$}\\
x &= 40 &&\justif{on divise par $12$}
\end{calculs}
\explique{$a=12>0$ : $-$ avant $40$, $+$ après.}{sur $\intff{0}{100}$ : première ligne de
$0$ à $100$, et non de $-\infty$ à $+\infty$}
\begin{center}
\begin{tikzpicture}[scale=0.85]
  \tkzTabInit[lgt=1.6,espcl=2.2]{$x$/1,$B(x)$/1}{$0$,$40$,$100$}
  \tkzTabLine{,-,z,+,}
\end{tikzpicture}
\end{center}

\cas{3)}
\explique{Elle gagne de l'argent quand $B(x)>0$, c'est-à-dire pour $x>40$.}{signe $+$ du
tableau}
\explique{Elle gagne de l'argent \res{à partir de $41$ bougies vendues}.}{$x$ est un nombre
entier ; pour $40$ bougies, $B(40)=0$ : elle ne gagne ni ne perd}
\end{exo}

\partie{Inéquations produit}

\begin{exo}[Signe d'un produit]
\explique{Pour chaque produit : le signe de chaque facteur (valeur qui l'annule, signe de
$a$), puis une ligne par facteur dans le tableau, et une dernière ligne pour le
produit.}{propriété 6 : même signe, produit positif ; signes contraires, produit négatif ;
un $0$ dans une colonne donne un $0$ au produit}

\cas{$A$}
\explique{$x-2=0$ pour $x=2$ ; $a=1>0$ : $-$ puis $+$.}{premier facteur}
\begin{calculs}
3x+6 &= 0 &&\justif{second facteur}\\
3x &= -6 &&\\
x &= -2 &&\justif{$a=3>0$ : $-$ puis $+$}
\end{calculs}
\explique{Première ligne : $-2$ puis $2$.}{valeurs rangées dans l'ordre croissant}
\begin{center}
\begin{tikzpicture}[scale=0.85]
  \tkzTabInit[lgt=3,espcl=2]{$x$/1,$x-2$/1,$3x+6$/1,$A(x)$/1}%
    {$-\infty$,$-2$,$2$,$+\infty$}
  \tkzTabLine{,-,t,-,z,+,}
  \tkzTabLine{,-,z,+,t,+,}
  \tkzTabLine{,+,z,-,z,+,}
\end{tikzpicture}
\end{center}
\explique{Dernière ligne, colonne par colonne : $(-)\times(-)$ donne $+$ ;
$(-)\times(+)$ donne $-$ ; $(+)\times(+)$ donne $+$.}{dans la ligne de $x-2$, on met un trait
sous $-2$ : ce facteur ne s'y annule pas}

\cas{$B$}
\begin{calculs}
4x-12 &= 0 &&\justif{premier facteur}\\
4x &= 12 &&\\
x &= 3 &&\justif{$a=4>0$ : $-$ puis $+$}
\end{calculs}
\explique{$-x+1=0$ pour $x=1$ ; $a=-1<0$ : $+$ puis $-$.}{second facteur}
\begin{center}
\begin{tikzpicture}[scale=0.85]
  \tkzTabInit[lgt=3,espcl=2]{$x$/1,$4x-12$/1,$-x+1$/1,$B(x)$/1}%
    {$-\infty$,$1$,$3$,$+\infty$}
  \tkzTabLine{,-,t,-,z,+,}
  \tkzTabLine{,+,z,-,t,-,}
  \tkzTabLine{,-,z,+,z,-,}
\end{tikzpicture}
\end{center}
\explique{$(-)\times(+)$ donne $-$ ; $(-)\times(-)$ donne $+$ ; $(+)\times(-)$ donne
$-$.}{contrôle : $B(0)=-12\times 1$, négatif}

\cas{$C$}
\explique{$5-x=0$ pour $x=5$ ; $a=-1<0$ : $+$ puis $-$.}{premier facteur, $5-x=-x+5$}
\begin{calculs}
2x+8 &= 0 &&\justif{second facteur}\\
2x &= -8 &&\\
x &= -4 &&\justif{$a=2>0$ : $-$ puis $+$}
\end{calculs}
\begin{center}
\begin{tikzpicture}[scale=0.85]
  \tkzTabInit[lgt=3,espcl=2]{$x$/1,$5-x$/1,$2x+8$/1,$C(x)$/1}%
    {$-\infty$,$-4$,$5$,$+\infty$}
  \tkzTabLine{,+,t,+,z,-,}
  \tkzTabLine{,-,z,+,t,+,}
  \tkzTabLine{,-,z,+,z,-,}
\end{tikzpicture}
\end{center}
\explique{$(+)\times(-)$ donne $-$ ; $(+)\times(+)$ donne $+$ ; $(-)\times(+)$ donne
$-$.}{contrôle : $C(0)=5\times 8$, positif}
\rem{Erreur fréquente : écrire les valeurs dans l'ordre où on les trouve ($5$ puis $-4$ pour
$C$). La première ligne est toujours rangée dans l'ordre croissant.}
\end{exo}

\begin{exo}[Résoudre]
\explique{Pour chaque inéquation : le tableau de signes du produit, puis on lit les signes
demandés.}{méthode 13 ; « $<0$ » : signes $-$ ; « $\ge 0$ » : signes $+$ et les $0$}

\cas{a)}
\explique{$x-4$ s'annule en $4$, $x+1$ en $-1$ ; tous deux : $-$ puis $+$.}{coefficients de
$x$ égaux à $1>0$}
\begin{center}
\begin{tikzpicture}[scale=0.85]
  \tkzTabInit[lgt=3.4,espcl=2]{$x$/1,$x-4$/1,$x+1$/1,$(x-4)(x+1)$/1}%
    {$-\infty$,$-1$,$4$,$+\infty$}
  \tkzTabLine{,-,t,-,z,+,}
  \tkzTabLine{,-,z,+,t,+,}
  \tkzTabLine{,+,z,-,z,+,}
\end{tikzpicture}
\end{center}
\explique{$\res{S=\intoo{-1}{4}}$.}{signes $-$ ; inégalité stricte : $-1$ et $4$ exclus}

\cas{b)}
\explique{$2x+6=0$ pour $x=-3$ ; $a=2>0$ : $-$ puis $+$. $x-5$ s'annule en $5$ : $-$ puis
$+$.}{$2x=-6$}
\begin{center}
\begin{tikzpicture}[scale=0.85]
  \tkzTabInit[lgt=3.4,espcl=2]{$x$/1,$2x+6$/1,$x-5$/1,$(2x+6)(x-5)$/1}%
    {$-\infty$,$-3$,$5$,$+\infty$}
  \tkzTabLine{,-,z,+,t,+,}
  \tkzTabLine{,-,t,-,z,+,}
  \tkzTabLine{,+,z,-,z,+,}
\end{tikzpicture}
\end{center}
\explique{$\res{S=\intof{-\infty}{-3}\cup\intfo{5}{+\infty}}$.}{signes $+$ et les $0$ :
« $\ge$ », $-3$ et $5$ compris}

\cas{c)}
\explique{$3-x$ s'annule en $3$ ; $a=-1<0$ : $+$ puis $-$. $x+2$ s'annule en $-2$ : $-$ puis
$+$.}{}
\begin{center}
\begin{tikzpicture}[scale=0.85]
  \tkzTabInit[lgt=3.4,espcl=2]{$x$/1,$3-x$/1,$x+2$/1,$(3-x)(x+2)$/1}%
    {$-\infty$,$-2$,$3$,$+\infty$}
  \tkzTabLine{,+,t,+,z,-,}
  \tkzTabLine{,-,z,+,t,+,}
  \tkzTabLine{,-,z,+,z,-,}
\end{tikzpicture}
\end{center}
\explique{$\res{S=\intoo{-2}{3}}$.}{signes $+$ ; inégalité stricte}

\cas{d)}
\explique{$-2x+4=0$ pour $x=2$ ; $a=-2<0$ : $+$ puis $-$. $3x-9=0$ pour $x=3$ ; $a=3>0$ :
$-$ puis $+$.}{$-2x=-4$ et $3x=9$}
\begin{center}
\begin{tikzpicture}[scale=0.85]
  \tkzTabInit[lgt=3.8,espcl=2]{$x$/1,$-2x+4$/1,$3x-9$/1,$(-2x+4)(3x-9)$/1}%
    {$-\infty$,$2$,$3$,$+\infty$}
  \tkzTabLine{,+,z,-,t,-,}
  \tkzTabLine{,-,t,-,z,+,}
  \tkzTabLine{,-,z,+,z,-,}
\end{tikzpicture}
\end{center}
\explique{$\res{S=\intof{-\infty}{2}\cup\intfo{3}{+\infty}}$.}{signes $-$ et les $0$ :
« $\le$ », $2$ et $3$ compris}
\rem{Avec « $\le$ » ou « $\ge$ », les valeurs où le produit s'annule sont solutions :
crochets fermés de leur côté. Du côté de l'infini, le crochet est toujours ouvert.}
\end{exo}

\begin{exo}[Trouver l'erreur]
\explique{Lina a mal rangé les valeurs de la première ligne : \res{$-1$ doit être écrit avant
$2$}, car $-1<2$.}{la première ligne est toujours dans l'ordre croissant}
\explique{Dans chaque ligne, le $0$ va sous la valeur qui annule le facteur ; puis on refait
la règle des signes.}{$x+1=0$ pour $x=-1$ ($-$ puis $+$) ; $2-x=0$ pour $x=2$ ($+$ puis
$-$)}
\begin{center}
\begin{tikzpicture}[scale=0.85]
  \tkzTabInit[lgt=3,espcl=2]{$x$/1,$x+1$/1,$2-x$/1,$(x+1)(2-x)$/1}%
    {$-\infty$,$-1$,$2$,$+\infty$}
  \tkzTabLine{,-,z,+,t,+,}
  \tkzTabLine{,+,t,+,z,-,}
  \tkzTabLine{,-,z,+,z,-,}
\end{tikzpicture}
\end{center}
\rem{Contrôle avec $x=0$, entre $-1$ et $2$ : $(0+1)(2-0)=2$, positif. Dans le tableau de
Lina, la colonne entre $2$ et $-1$ n'a pas de sens : il n'y a aucun nombre plus grand que $2$
et plus petit que $-1$.}
\end{exo}

\partie{Inéquations quotient}

\begin{exo}[Signe d'un quotient]
\explique{Pour chaque quotient : la valeur interdite annule le dénominateur ; elle est
marquée d'une double barre dans la ligne du quotient.}{propriété 7 : même règle des signes
qu'un produit, mais on ne divise pas par $0$}

\cas{$A$}
\explique{Valeur interdite : $x-2=0$ pour $\res{x=2}$.}{dénominateur nul}
\explique{$x+3=0$ pour $x=-3$ : $-$ puis $+$. $x-2$ : $-$ puis $+$.}{coefficients de $x$
égaux à $1>0$}
\begin{center}
\begin{tikzpicture}[scale=0.85]
  \tkzTabInit[lgt=2.6,espcl=2]{$x$/1,$x+3$/1,$x-2$/1,$A(x)$/1}%
    {$-\infty$,$-3$,$2$,$+\infty$}
  \tkzTabLine{,-,z,+,t,+,}
  \tkzTabLine{,-,t,-,z,+,}
  \tkzTabLine{,+,z,-,d,+,}
\end{tikzpicture}
\end{center}
\explique{En $-3$, le numérateur est nul : $A(-3)=0$. En $2$, double barre : $A(2)$
n'existe pas.}{$0$ au quotient sous le zéro du numérateur ; double barre sous celui du
dénominateur}

\cas{$B$}
\explique{Valeur interdite : $x+4=0$ pour $\res{x=-4}$.}{dénominateur nul}
\begin{calculs}
6-2x &= 0 &&\justif{numérateur}\\
-2x &= -6 &&\\
x &= 3 &&\justif{$a=-2<0$ : $+$ puis $-$}
\end{calculs}
\begin{center}
\begin{tikzpicture}[scale=0.85]
  \tkzTabInit[lgt=2.6,espcl=2]{$x$/1,$6-2x$/1,$x+4$/1,$B(x)$/1}%
    {$-\infty$,$-4$,$3$,$+\infty$}
  \tkzTabLine{,+,t,+,z,-,}
  \tkzTabLine{,-,z,+,t,+,}
  \tkzTabLine{,-,d,+,z,-,}
\end{tikzpicture}
\end{center}

\cas{$C$}
\explique{Valeur interdite : $5-x=0$ pour $\res{x=5}$.}{dénominateur nul}
\explique{$3x=0$ pour $x=0$ ; $a=3>0$ : $-$ puis $+$. $5-x$ : $a=-1<0$, $+$ puis $-$.}{}
\begin{center}
\begin{tikzpicture}[scale=0.85]
  \tkzTabInit[lgt=2.6,espcl=2]{$x$/1,$3x$/1,$5-x$/1,$C(x)$/1}%
    {$-\infty$,$0$,$5$,$+\infty$}
  \tkzTabLine{,-,z,+,t,+,}
  \tkzTabLine{,+,t,+,z,-,}
  \tkzTabLine{,-,z,+,d,-,}
\end{tikzpicture}
\end{center}
\rem{Contrôle avec une valeur de chaque colonne, par exemple pour $C$ : $C(-5)=\dfrac{-15}{10}$,
négatif ; $C(1)=\dfrac{3}{4}$, positif ; $C(10)=\dfrac{30}{-5}$, négatif.}
\end{exo}

\begin{exo}[Résoudre]
\explique{Valeur interdite, tableau de signes du quotient, puis lecture : la valeur interdite
n'est jamais solution.}{méthode 14 du cours}

\cas{a)}
\explique{Valeur interdite : $x=-3$. $x-1$ s'annule en $1$ : $-$ puis $+$ ; $x+3$ : $-$ puis
$+$.}{}
\begin{center}
\begin{tikzpicture}[scale=0.85]
  \tkzTabInit[lgt=2.6,espcl=2]{$x$/1,$x-1$/1,$x+3$/1,$\dfrac{x-1}{x+3}$/1.6}%
    {$-\infty$,$-3$,$1$,$+\infty$}
  \tkzTabLine{,-,t,-,z,+,}
  \tkzTabLine{,-,z,+,t,+,}
  \tkzTabLine{,+,d,-,z,+,}
\end{tikzpicture}
\end{center}
\explique{$\res{S=\intoo{-\infty}{-3}\cup\intfo{1}{+\infty}}$.}{signes $+$ ; $1$ compris
($\ge$) ; $-3$ exclu (interdite)}

\cas{b)}
\explique{Valeur interdite : $x=6$. $2x+8=0$ pour $x=-4$ ; $a=2>0$ : $-$ puis $+$.}{$2x=-8$}
\begin{center}
\begin{tikzpicture}[scale=0.85]
  \tkzTabInit[lgt=2.6,espcl=2]{$x$/1,$2x+8$/1,$x-6$/1,$\dfrac{2x+8}{x-6}$/1.6}%
    {$-\infty$,$-4$,$6$,$+\infty$}
  \tkzTabLine{,-,z,+,t,+,}
  \tkzTabLine{,-,t,-,z,+,}
  \tkzTabLine{,+,z,-,d,+,}
\end{tikzpicture}
\end{center}
\explique{$\res{S=\intoo{-4}{6}}$.}{signes $-$ ; inégalité stricte : $-4$ exclu ; $6$
interdit}

\cas{c)}
\explique{Valeur interdite : $2x-2=0$ pour $x=1$. $4-x$ s'annule en $4$ ; $a=-1<0$ : $+$ puis
$-$.}{$2x-2$ : $a=2>0$, $-$ puis $+$}
\begin{center}
\begin{tikzpicture}[scale=0.85]
  \tkzTabInit[lgt=2.6,espcl=2]{$x$/1,$4-x$/1,$2x-2$/1,$\dfrac{4-x}{2x-2}$/1.6}%
    {$-\infty$,$1$,$4$,$+\infty$}
  \tkzTabLine{,+,t,+,z,-,}
  \tkzTabLine{,-,z,+,t,+,}
  \tkzTabLine{,-,d,+,z,-,}
\end{tikzpicture}
\end{center}
\explique{$\res{S=\intoo{-\infty}{1}\cup\intfo{4}{+\infty}}$.}{signes $-$ ; $4$ compris
($\le$) ; $1$ exclu (interdite)}

\cas{d)}
\explique{Valeur interdite : $x=-5$.}{$x+5=0$}
\begin{calculs}
-3x-6 &= 0 &&\justif{numérateur}\\
-3x &= 6 &&\justif{on ajoute $6$}\\
x &= -2 &&\justif{on divise par $-3$ ; $a=-3<0$ : $+$ puis $-$}
\end{calculs}
\begin{center}
\begin{tikzpicture}[scale=0.85]
  \tkzTabInit[lgt=2.6,espcl=2]{$x$/1,$-3x-6$/1,$x+5$/1,$\dfrac{-3x-6}{x+5}$/1.6}%
    {$-\infty$,$-5$,$-2$,$+\infty$}
  \tkzTabLine{,+,t,+,z,-,}
  \tkzTabLine{,-,z,+,t,+,}
  \tkzTabLine{,-,d,+,z,-,}
\end{tikzpicture}
\end{center}
\explique{$\res{S=\intoo{-5}{-2}}$.}{signes $+$ ; inégalité stricte}
\rem{La valeur interdite est toujours exclue : crochet ouvert de son côté, même avec
« $\le$ » ou « $\ge$ ».}
\end{exo}

\begin{exo}[Produit, puis quotient]
\cas{1)}
\explique{$3x+6=0$ pour $x=-2$ ; $a=3>0$ : $-$ puis $+$.}{$3x=-6$}
\explique{$-x+4=0$ pour $x=4$ ; $a=-1<0$ : $+$ puis $-$.}{}
\begin{center}
\begin{tikzpicture}[scale=0.85]
  \tkzTabInit[lgt=3.6,espcl=2]{$x$/1,$3x+6$/1,$-x+4$/1,$(3x+6)(-x+4)$/1}%
    {$-\infty$,$-2$,$4$,$+\infty$}
  \tkzTabLine{,-,z,+,t,+,}
  \tkzTabLine{,+,t,+,z,-,}
  \tkzTabLine{,-,z,+,z,-,}
\end{tikzpicture}
\end{center}
\explique{$\res{S=\intoo{-2}{4}}$.}{signes $+$ ; inégalité stricte}

\cas{2)}
\explique{Le quotient n'existe pas pour $\res{x=4}$.}{le dénominateur $-x+4$ y est nul}

\cas{3)}
\explique{Mêmes facteurs, donc mêmes signes que le produit, sauf en $4$ : double barre au
lieu de $0$.}{on ajoute une ligne au tableau de 1)}
\begin{center}
\begin{tikzpicture}[scale=0.85]
  \tkzTabInit[lgt=3.6,espcl=2]{$x$/1,$3x+6$/1,$-x+4$/1,$(3x+6)(-x+4)$/1,$\dfrac{3x+6}{-x+4}$/1.6}%
    {$-\infty$,$-2$,$4$,$+\infty$}
  \tkzTabLine{,-,z,+,t,+,}
  \tkzTabLine{,+,t,+,z,-,}
  \tkzTabLine{,-,z,+,z,-,}
  \tkzTabLine{,-,z,+,d,-,}
\end{tikzpicture}
\end{center}
\explique{$\res{S=\intof{-\infty}{-2}\cup\intoo{4}{+\infty}}$.}{signes $-$ et le $0$ en
$-2$ (« $\le$ ») ; $4$ exclu (interdite)}
\rem{Un produit et un quotient des mêmes facteurs ont le même signe ; seule change la valeur
qui annule le dénominateur.}
\end{exo}

\partie{Inéquations avancées}

\begin{exo}[Factoriser d'abord]
\explique{Le tableau de signes compare un produit à $0$ : on passe tout à gauche, puis on
factorise par le facteur commun.}{méthode 15 du cours}

\cas{a)}
\begin{calculs}
x^{2} &\ge 3x &&\justif{on recopie l'inéquation}\\
x^{2}-3x &\ge 0 &&\justif{on passe $3x$ à gauche}\\
x(x-3) &\ge 0 &&\justif{facteur commun $x$}
\end{calculs}
\explique{$x$ s'annule en $0$, $x-3$ en $3$ ; tous deux : $-$ puis $+$.}{coefficients de $x$
positifs}
\begin{center}
\begin{tikzpicture}[scale=0.85]
  \tkzTabInit[lgt=2.6,espcl=2]{$x$/1,$x$/1,$x-3$/1,$x(x-3)$/1}%
    {$-\infty$,$0$,$3$,$+\infty$}
  \tkzTabLine{,-,z,+,t,+,}
  \tkzTabLine{,-,t,-,z,+,}
  \tkzTabLine{,+,z,-,z,+,}
\end{tikzpicture}
\end{center}
\explique{$\res{S=\intof{-\infty}{0}\cup\intfo{3}{+\infty}}$.}{signes $+$ et les $0$}

\cas{b)}
\begin{calculs}
2x^{2} &< 8x &&\justif{on recopie l'inéquation}\\
2x^{2}-8x &< 0 &&\justif{on passe $8x$ à gauche}\\
2x(x-4) &< 0 &&\justif{facteur commun $2x$}
\end{calculs}
\explique{$2x$ s'annule en $0$, $x-4$ en $4$ ; tous deux : $-$ puis $+$.}{}
\begin{center}
\begin{tikzpicture}[scale=0.85]
  \tkzTabInit[lgt=2.6,espcl=2]{$x$/1,$2x$/1,$x-4$/1,$2x(x-4)$/1}%
    {$-\infty$,$0$,$4$,$+\infty$}
  \tkzTabLine{,-,z,+,t,+,}
  \tkzTabLine{,-,t,-,z,+,}
  \tkzTabLine{,+,z,-,z,+,}
\end{tikzpicture}
\end{center}
\explique{$\res{S=\intoo{0}{4}}$.}{signes $-$ ; inégalité stricte}

\cas{c)}
\begin{calculs}
x^{2} &\le -5x &&\justif{on recopie l'inéquation}\\
x^{2}+5x &\le 0 &&\justif{on ajoute $5x$ des deux côtés}\\
x(x+5) &\le 0 &&\justif{facteur commun $x$}
\end{calculs}
\explique{$x+5$ s'annule en $-5$, $x$ en $0$ ; tous deux : $-$ puis $+$.}{valeurs rangées :
$-5$ puis $0$}
\begin{center}
\begin{tikzpicture}[scale=0.85]
  \tkzTabInit[lgt=2.6,espcl=2]{$x$/1,$x$/1,$x+5$/1,$x(x+5)$/1}%
    {$-\infty$,$-5$,$0$,$+\infty$}
  \tkzTabLine{,-,t,-,z,+,}
  \tkzTabLine{,-,z,+,t,+,}
  \tkzTabLine{,+,z,-,z,+,}
\end{tikzpicture}
\end{center}
\explique{$\res{S=\intff{-5}{0}}$.}{signes $-$ et les $0$}
\rem{Ne jamais diviser par $x$ : on ne connaît pas son signe (le sens changerait-il ?), et
on perdrait la solution $x=0$. Par exemple, en a), $x=0$ vérifie bien $0\ge 0$.}
\end{exo}

\begin{exo}[Avec une identité remarquable]
\explique{On passe tout à gauche pour comparer à $0$, puis on factorise avec
$a^{2}-b^{2}=(a-b)(a+b)$.}{méthode 15 du cours}

\cas{a)}
\begin{calculs}
x^{2} &< 25 &&\justif{on recopie l'inéquation}\\
x^{2}-25 &< 0 &&\justif{on passe $25$ à gauche}\\
(x-5)(x+5) &< 0 &&\justif{$a^{2}-b^{2}$ avec $a=x$ et $b=5$}
\end{calculs}
\explique{$x+5$ s'annule en $-5$, $x-5$ en $5$ ; tous deux : $-$ puis $+$.}{}
\begin{center}
\begin{tikzpicture}[scale=0.85]
  \tkzTabInit[lgt=3,espcl=2]{$x$/1,$x-5$/1,$x+5$/1,$(x-5)(x+5)$/1}%
    {$-\infty$,$-5$,$5$,$+\infty$}
  \tkzTabLine{,-,t,-,z,+,}
  \tkzTabLine{,-,z,+,t,+,}
  \tkzTabLine{,+,z,-,z,+,}
\end{tikzpicture}
\end{center}
\explique{$\res{S=\intoo{-5}{5}}$.}{signes $-$ ; inégalité stricte}
\rem{Erreur fréquente : répondre $x<5$. Mais $x=-6$ ne convient pas : $(-6)^{2}=36$, et
$36>25$.}

\cas{b)}
\begin{calculs}
(x-3)^{2} &\ge 4 &&\justif{on recopie l'inéquation}\\
(x-3)^{2}-4 &\ge 0 &&\justif{on passe $4$ à gauche}\\
\big[(x-3)-2\big]\big[(x-3)+2\big] &\ge 0 &&\justif{$a^{2}-b^{2}$ avec $a=x-3$ et $b=2$}\\
(x-5)(x-1) &\ge 0 &&\justif{on réduit chaque crochet}
\end{calculs}
\explique{$x-1$ s'annule en $1$, $x-5$ en $5$ ; tous deux : $-$ puis $+$.}{}
\begin{center}
\begin{tikzpicture}[scale=0.85]
  \tkzTabInit[lgt=3,espcl=2]{$x$/1,$x-5$/1,$x-1$/1,$(x-5)(x-1)$/1}%
    {$-\infty$,$1$,$5$,$+\infty$}
  \tkzTabLine{,-,t,-,z,+,}
  \tkzTabLine{,-,z,+,t,+,}
  \tkzTabLine{,+,z,-,z,+,}
\end{tikzpicture}
\end{center}
\explique{$\res{S=\intof{-\infty}{1}\cup\intfo{5}{+\infty}}$.}{signes $+$ et les $0$}

\cas{c)}
\begin{calculs}
(2x+1)^{2} &\le 9 &&\justif{on recopie l'inéquation}\\
(2x+1)^{2}-9 &\le 0 &&\justif{on passe $9$ à gauche}\\
\big[(2x+1)-3\big]\big[(2x+1)+3\big] &\le 0 &&\justif{$a^{2}-b^{2}$ avec $a=2x+1$ et $b=3$}\\
(2x-2)(2x+4) &\le 0 &&\justif{on réduit chaque crochet}
\end{calculs}
\explique{$2x-2=0$ pour $x=1$ ; $2x+4=0$ pour $x=-2$ ; tous deux : $-$ puis $+$.}{$a=2>0$
pour les deux facteurs}
\begin{center}
\begin{tikzpicture}[scale=0.85]
  \tkzTabInit[lgt=3,espcl=2]{$x$/1,$2x-2$/1,$2x+4$/1,$(2x-2)(2x+4)$/1}%
    {$-\infty$,$-2$,$1$,$+\infty$}
  \tkzTabLine{,-,t,-,z,+,}
  \tkzTabLine{,-,z,+,t,+,}
  \tkzTabLine{,+,z,-,z,+,}
\end{tikzpicture}
\end{center}
\explique{$\res{S=\intff{-2}{1}}$.}{signes $-$ et les $0$ ; contrôle avec $x=0$ : $1\le 9$
est vrai}
\end{exo}

\begin{exo}[Comparer un quotient à un nombre]
\explique{On passe le nombre à gauche, puis on met au même dénominateur pour obtenir un seul
quotient, comparé à $0$.}{méthode 15 du cours}

\cas{a)}
\explique{Valeur interdite : $x=1$.}{$x-1=0$}
\begin{calculs}
\frac{x+5}{x-1}-2 &\ge 0 &&\justif{on passe $2$ à gauche}\\
\frac{x+5}{x-1}-\frac{2\fois{(x-1)}}{1\fois{(x-1)}} &\ge 0 &&\justif{même dénominateur}\\
\frac{x+5-2x+2}{x-1} &\ge 0 &&\justif{attention au signe $-$ devant la fraction :
$-2(x-1)=-2x+2$}\\
\frac{-x+7}{x-1} &\ge 0 &&\justif{on réduit}
\end{calculs}
\explique{$-x+7$ s'annule en $7$ ; $a=-1<0$ : $+$ puis $-$. $x-1$ : $-$ puis $+$.}{}
\begin{center}
\begin{tikzpicture}[scale=0.85]
  \tkzTabInit[lgt=2.6,espcl=2]{$x$/1,$-x+7$/1,$x-1$/1,$\dfrac{-x+7}{x-1}$/1.6}%
    {$-\infty$,$1$,$7$,$+\infty$}
  \tkzTabLine{,+,t,+,z,-,}
  \tkzTabLine{,-,z,+,t,+,}
  \tkzTabLine{,-,d,+,z,-,}
\end{tikzpicture}
\end{center}
\explique{$\res{S=\intof{1}{7}}$.}{signes $+$ ; $7$ compris ($\ge$) ; $1$ exclu
(interdite)}
\rem{Contrôle avec $x=3$ : $\dfrac{8}{2}=4$, et $4\ge 2$ ; $3$ est bien dans $S$.}

\cas{b)}
\explique{Valeur interdite : $x=-2$.}{$x+2=0$}
\begin{calculs}
\frac{3}{x+2}-1 &< 0 &&\justif{on passe $1$ à gauche}\\
\frac{3}{x+2}-\frac{1\fois{(x+2)}}{1\fois{(x+2)}} &< 0 &&\justif{même dénominateur}\\
\frac{3-x-2}{x+2} &< 0 &&\justif{$-(x+2)=-x-2$}\\
\frac{1-x}{x+2} &< 0 &&\justif{on réduit}
\end{calculs}
\explique{$1-x$ s'annule en $1$ ; $a=-1<0$ : $+$ puis $-$. $x+2$ : $-$ puis $+$.}{}
\begin{center}
\begin{tikzpicture}[scale=0.85]
  \tkzTabInit[lgt=2.6,espcl=2]{$x$/1,$1-x$/1,$x+2$/1,$\dfrac{1-x}{x+2}$/1.6}%
    {$-\infty$,$-2$,$1$,$+\infty$}
  \tkzTabLine{,+,t,+,z,-,}
  \tkzTabLine{,-,z,+,t,+,}
  \tkzTabLine{,-,d,+,z,-,}
\end{tikzpicture}
\end{center}
\explique{$\res{S=\intoo{-\infty}{-2}\cup\intoo{1}{+\infty}}$.}{signes $-$ ; inégalité
stricte ; $-2$ interdit}
\rem{Erreur fréquente : multiplier les deux membres par $x+2$ pour « chasser » le
dénominateur. On ne connaît pas le signe de $x+2$ : on ne sait pas si le sens change.}
\end{exo}

\begin{exo}[Coût moyen]
\cas{1)}
\begin{calculs}
M(10) &= \frac{2\times 10+60}{10} &&\justif{on remplace $x$ par $10$}\\
M(10) &= \frac{80}{10} &&\\
M(10) &= 8 &&
\end{calculs}
\begin{calculs}
M(30) &= \frac{2\times 30+60}{30} &&\justif{on remplace $x$ par $30$}\\
M(30) &= \frac{120}{30} &&\\
M(30) &= 4 &&
\end{calculs}
\explique{Pour $10$ figurines, chacune revient à \res{$8$~€} ; pour $30$ figurines, à
\res{$4$~€}.}{les $60$~€ de réglages se partagent entre plus de figurines}

\cas{2)}
\begin{calculs}
\frac{2x+60}{x} &\le 5 &&\justif{on traduit $M(x)\le 5$}\\
\frac{2x+60}{x}-5 &\le 0 &&\justif{on passe $5$ à gauche}\\
\frac{2x+60}{x}-\frac{5\fois{x}}{1\fois{x}} &\le 0 &&\justif{même dénominateur}\\
\frac{2x+60-5x}{x} &\le 0 &&\\
\frac{60-3x}{x} &\le 0 &&\justif{on réduit : c'est l'inéquation annoncée}
\end{calculs}

\cas{3)}
\explique{Pour $x>0$, le dénominateur $x$ est positif.}{il n'y a pas de valeur interdite sur
$\intoo{0}{+\infty}$ ; $0$ est exclu}
\begin{calculs}
60-3x &= 0 &&\justif{numérateur}\\
-3x &= -60 &&\\
x &= 20 &&\justif{$a=-3<0$ : $+$ puis $-$}
\end{calculs}
\begin{center}
\begin{tikzpicture}[scale=0.85]
  \tkzTabInit[lgt=2.6,espcl=2.2]{$x$/1,$60-3x$/1,$x$/1,$\dfrac{60-3x}{x}$/1.6}%
    {$0$,$20$,$+\infty$}
  \tkzTabLine{,+,z,-,}
  \tkzTabLine{z,+,t,+,}
  \tkzTabLine{d,+,z,-,}
\end{tikzpicture}
\end{center}
\explique{$\dfrac{60-3x}{x}\le 0$ pour $x\ge 20$ : $S=\intfo{20}{+\infty}$.}{signes $-$ et le
$0$}
\explique{Le coût moyen est d'au plus $5$~€ \res{à partir de $20$ figurines}.}{contrôle :
$M(20)=\dfrac{100}{20}$, soit $5$ ; « au plus » : $20$ convient}
\end{exo}

\partie{Problèmes d'inéquations}

\begin{exo}[Moyenne trimestrielle]
\cas{1)}
\explique{$\res{M(n)=\dfrac{9+12+10+n}{4}}$, soit $\res{M(n)=\dfrac{31+n}{4}}$.}{somme des
quatre notes, divisée par $4$}

\cas{2)}
\explique{« Au moins $12$ » se traduit par $M(n)\ge 12$.}{au moins : $12$ ou plus}
\begin{calculs}
\frac{31+n}{4} &\ge 12 &&\justif{on recopie l'inéquation}\\
31+n &\ge 48 &&\justif{on multiplie par $4>0$ : même sens}\\
n &\ge 17 &&\justif{on enlève $31$}
\end{calculs}
\explique{Une note est comprise entre $0$ et $20$ : il lui faut \res{une note entre $17$ et
$20$}.}{$S=\intff{17}{20}$ ; contrôle : $M(17)=\dfrac{48}{4}$, soit $12$}

\cas{3)}
\begin{calculs}
\frac{31+n}{4} &\ge 15 &&\justif{on traduit $M(n)\ge 15$}\\
31+n &\ge 60 &&\justif{on multiplie par $4>0$}\\
n &\ge 29 &&\justif{on enlève $31$}
\end{calculs}
\explique{\res{Non} : il faudrait au moins $29$ sur $20$, ce qui est impossible.}{même avec
$20$, $M(20)=\dfrac{51}{4}$, soit $12{,}75$}
\end{exo}

\begin{exo}[Rectangle ou carré]
\cas{1)}
\begin{calculs}
P(x) &= 2(x+5)+2\times 3 &&\justif{deux longueurs et deux largeurs}\\
P(x) &= 2x+10+6 &&\justif{on développe}\\
P(x) &= 2x+16 &&
\end{calculs}
\explique{$\res{P(x)=2x+16}$ et $\res{Q(x)=4x}$.}{le carré a quatre côtés égaux}

\cas{2)}
\explique{« Le périmètre du carré est plus grand que celui du rectangle » : $Q(x)>P(x)$.}{on
traduit la phrase par une inéquation}
\begin{calculs}
4x &> 2x+16 &&\justif{on recopie l'inéquation}\\
4x-2x &> 16 &&\justif{les $x$ à gauche}\\
2x &> 16 &&\justif{on réduit}\\
x &> 8 &&\justif{on divise par $2>0$ : même sens}
\end{calculs}
\explique{Le périmètre du carré est plus grand \res{pour $x>8$}, soit
$S=\intoo{8}{+\infty}$.}{la condition $x>0$ est bien respectée}

\cas{3)}
\begin{calculs}
P(10) &= 2\times 10+16 &&\\
P(10) &= 36 &&
\end{calculs}
\begin{calculs}
Q(10) &= 4\times 10 &&\\
Q(10) &= 40 &&
\end{calculs}
\explique{$40>36$ : pour $x=10$, le carré a bien le plus grand périmètre, et $10>8$.}{la
vérification est cohérente avec 2)}
\end{exo}

\begin{exo}[Trottinettes en libre-service]
\cas{1)}
\explique{$\res{A(x)=0{,}2x+1}$.}{$0{,}20$~€ par minute, soit $0{,}2x$, plus $1$~€ de
déblocage}
\explique{$\res{B(x)=0{,}3x}$.}{$0{,}30$~€ par minute, rien d'autre : fonction linéaire}

\cas{2)}
\explique{$5$ minutes : $A(5)=1+1$, soit \res{$2$~€} ; $B(5)=$ \res{$1{,}50$~€}.}{$0{,}2\times
5=1$ et $0{,}3\times 5=1{,}5$}
\explique{$20$ minutes : $A(20)=4+1$, soit \res{$5$~€} ; $B(20)=$ \res{$6$~€}.}{$0{,}2\times
20=4$ et $0{,}3\times 20=6$}
\explique{Pour $5$ minutes, B est moins cher ; pour $20$ minutes, c'est A.}{le classement
s'inverse entre les deux : on cherche où}

\cas{3)}
\begin{calculs}
0{,}2x+1 &< 0{,}3x &&\justif{on traduit $A(x)<B(x)$}\\
0{,}2x-0{,}3x &< -1 &&\justif{les $x$ à gauche, les nombres à droite}\\
-0{,}1x &< -1 &&\justif{on réduit}\\
x &> 10 &&\justif{on divise par $-0{,}1<0$ : le sens change}
\end{calculs}
\explique{Avec $0\le x\le 30$ : $\res{S=\intof{10}{30}}$.}{on garde les solutions de
l'intervalle de l'énoncé}

\cas{4)}
\explique{Pour un trajet de \res{moins de $10$ minutes}, on choisit \res{B} ; pour \res{plus de
$10$ minutes}, on choisit \res{A}.}{d'après 3), A est moins cher pour $x>10$}
\explique{Pour $10$ minutes exactement, les deux coûtent $3$~€.}{$A(10)=2+1$ et
$B(10)=3$}
\end{exo}

\partiesansnum{Problème de synthèse}

\begin{exo}[Étude complète d'une fonction]
\cas{1)}
\explique{On calcule les images de $0$ à $6$, en avançant de $1$ (le pas).}{sept valeurs}
\begin{calculs}
f(0) &= (0-1)(0-5) &&\justif{on remplace $x$ par $0$}\\
f(0) &= (-1)\times(-5) = 5 &&\justif{produit de deux négatifs : positif}
\end{calculs}
\begin{calculs}
f(3) &= (3-1)(3-5) &&\justif{on remplace $x$ par $3$}\\
f(3) &= 2\times(-2) = -4 &&
\end{calculs}
\begin{center}
\renewcommand{\arraystretch}{1.3}
\begin{tabular}{|c|*{7}{c|}}
\hline
$x$ & $0$ & $1$ & $2$ & $3$ & $4$ & $5$ & $6$\\ \hline
$f(x)$ & $5$ & $0$ & $-3$ & $-4$ & $-3$ & $0$ & $5$\\ \hline
\end{tabular}
\end{center}
\explique{En $1$ et en $5$, un facteur est nul, donc le produit aussi.}{$f(1)=0\times(-4)$ et
$f(5)=4\times 0$}

\cas{2)}
\begin{lecture}
On place les sept points du tableau, puis on les relie par une courbe sans angle.

\smallskip
\gris{tangente horizontale au point le plus bas, $(3\,;-4)$}
\end{lecture}
\begin{figlecture}
\begin{tikzpicture}[x=0.36cm,y=0.36cm]
  \repere{-1}{-5}{7}{6}
  \gradx{1,2,3,4,5,6} \grady{-4,-3,-2,-1,1,2,3,4,5}
  \draw[coulCourbe,line width=1pt,domain=0:6,samples=60]
    plot (\x,{(\x-1)*(\x-5)});
  \foreach \a/\b in {0/5,1/0,2/-3,3/-4,4/-3,5/0,6/5}
    \path[coulCourbe] (\a,\b) node[croix]{};
  \node[coulCourbe,font=\footnotesize,right] at (6,5) {$\Cf$};
\end{tikzpicture}
\end{figlecture}

\cas{3)}
\explique{De $x=0$ à $x=3$, la courbe descend : $f$ est \res{décroissante sur
$\intff{0}{3}$}.}{lecture sur la courbe de 2)}
\explique{De $x=3$ à $x=6$, elle monte : $f$ est \res{croissante sur $\intff{3}{6}$}.}{}
\explique{Tableau de variations :}{bornes $0$ et $6$ ; changement de sens en $3$ ; images
lues dans le tableau de valeurs}
\begin{center}
\begin{tikzpicture}[scale=0.85]
  \tkzTabInit[lgt=1.4,espcl=2.2,deltacl=0.6]{$x$/1,$f(x)$/2}{$0$,$3$,$6$}
  \tkzTabVar{+/$5$,-/$-4$,+/$5$}
\end{tikzpicture}
\end{center}

\cas{4)}
\explique{Le maximum de $f$ est \res{$5$, atteint en $0$ et en $6$}.}{plus grand nombre de la
seconde ligne : il est atteint deux fois, aux deux bords}
\explique{Le minimum de $f$ est \res{$-4$, atteint en $3$}.}{plus petit nombre de la seconde
ligne : le point le plus bas de la courbe}

\cas{5)}
\explique{$x-1$ s'annule en $1$, $x-5$ en $5$ ; tous deux : $-$ puis $+$.}{coefficients de
$x$ égaux à $1>0$ ; première ligne de $0$ à $6$}
\begin{center}
\begin{tikzpicture}[scale=0.85]
  \tkzTabInit[lgt=2.6,espcl=2]{$x$/1,$x-1$/1,$x-5$/1,$f(x)$/1}%
    {$0$,$1$,$5$,$6$}
  \tkzTabLine{,-,z,+,t,+,}
  \tkzTabLine{,-,t,-,z,+,}
  \tkzTabLine{,+,z,-,z,+,}
\end{tikzpicture}
\end{center}

\cas{6)}
\begin{lecture}
$f(x)<0$ : on lit les signes $-$ du tableau de 5).

\smallskip
$\res{S=\intoo{1}{5}}$.

\smallskip
\gris{inégalité stricte : $1$ et $5$ exclus}

\medskip
Sur la courbe : elle est au-dessous de l'axe des abscisses entre $x=1$ et $x=5$ (en rouge
sur l'axe), ce qui confirme le calcul.
\end{lecture}
\begin{figlecture}
\begin{tikzpicture}[x=0.36cm,y=0.36cm]
  \repere{-1}{-5}{7}{6}
  \gradx{1,2,3,4,5,6} \grady{-4,-3,-2,-1,1,2,3,4,5}
  \draw[coulCourbe,line width=1pt,domain=0:6,samples=60]
    plot (\x,{(\x-1)*(\x-5)});
  \morceau{1}{5}\crochet{1}{]}\crochet{5}{[}
\end{tikzpicture}
\end{figlecture}
\rem{Le tableau de variations (question 3) et le tableau de signes (question 5) ne
contiennent pas les mêmes valeurs de $x$ : $3$ pour les variations, $1$ et $5$ pour le
signe.}
\end{exo}

\end{document}
